\documentclass[11pt]{article}
\usepackage[a4paper,margin=2.8cm]{geometry}
\usepackage{amsmath,amssymb,amsthm,mathtools}
\usepackage{booktabs,array,enumitem}
\usepackage{microtype}
\usepackage[skip=0.75\baselineskip plus 2pt minus 1pt]{parskip}
\usepackage{titlesec}
\usepackage[font=small,labelfont=bf,margin=1em]{caption}
\usepackage[hidelinks]{hyperref}
\hypersetup{pdftitle={Neighbourhood counting and new bounds for small three-colour Ramsey numbers},
  pdfauthor={Lachlan Dauth}, pdfsubject={Mathematics - Combinatorics (05C55)},
  pdfkeywords={Multicolour Ramsey numbers, neighbourhood recurrence, degree profiles, flag algebras, semidefinite programming, computer-assisted proof}}

% Layout: block paragraphs with space between them, roomier lists,
% theorems and section headings, slightly opened line spacing.
\linespread{1.06}
\setlist{itemsep=0.35em,topsep=0.5em,parsep=0pt,partopsep=0pt}
\setlist[enumerate]{leftmargin=2.2em}
\setlist[itemize]{leftmargin=1.8em}
\titlespacing*{\section}{0pt}{2.6ex plus 1ex minus .2ex}{1.4ex plus .2ex}
\titlespacing*{\subsection}{0pt}{2.0ex plus 1ex minus .2ex}{1.0ex plus .2ex}
\titlespacing*{\paragraph}{0pt}{1.2ex plus .5ex minus .2ex}{0.8em}
\renewcommand{\arraystretch}{1.15}
% abstract: block paragraphs with a gap, like the body
\renewenvironment{abstract}{\small\begin{center}\bfseries\abstractname\end{center}%
  \list{}{\setlength{\leftmargin}{2.5em}\setlength{\rightmargin}{\leftmargin}%
  \setlength{\listparindent}{0pt}\setlength{\parsep}{0.6em plus 1pt}\setlength{\topsep}{0pt}}\item\relax}{\endlist}
\newtheoremstyle{spacedplain}{1.1em}{1.1em}{\itshape}{}{\bfseries}{.}{0.6em}{}
\newtheoremstyle{spaceddef}{1.1em}{1.1em}{\normalfont}{}{\bfseries}{.}{0.6em}{}
\newtheoremstyle{spacedremark}{1.0em}{1.0em}{\normalfont}{}{\itshape}{.}{0.6em}{}
\theoremstyle{spacedplain}

\newtheorem{theorem}{Theorem}[section]
\newtheorem{lemma}[theorem]{Lemma}
\newtheorem{corollary}[theorem]{Corollary}
\newtheorem{proposition}[theorem]{Proposition}
\theoremstyle{spaceddef}
\newtheorem{definition}[theorem]{Definition}
\theoremstyle{spacedremark}
\newtheorem{remark}[theorem]{Remark}

\newcommand{\R}{\mathrm{R}}
\newcommand{\bt}{\mathbf{t}}
\newcommand{\e}{\mathbf{e}}
\newcommand{\Nr}{N_{1}}
\newcommand{\Nb}{N_{2}}
\newcommand{\Ng}{N_{3}}
% Review copy (wide notes margin, line numbers): compile the _markup.tex wrapper
\ifdefined\markupcopy% Layout for the printed review copy. Loaded only when \markupcopy is defined.
% Content, numbering and wording are identical to the clean draft.
\geometry{a4paper,left=2.0cm,right=6.2cm,top=2.6cm,bottom=2.6cm,headheight=14pt,headsep=0.6cm}
\usepackage{xcolor}
\usepackage{setspace}
\setstretch{1.7}
\usepackage{etoolbox}
\AtBeginEnvironment{tabular}{\setstretch{1.15}}
\AtBeginEnvironment{thebibliography}{\setstretch{1.25}}
% wide tables may not spill into the notes column
\usepackage{adjustbox}
\let\mkoldtabular\tabular
\let\mkoldendtabular\endtabular
\renewenvironment{tabular}{\begin{adjustbox}{max width=\linewidth}\mkoldtabular}{\mkoldendtabular\end{adjustbox}}
% line numbers on every line, for referencing comments
\usepackage[mathlines]{lineno}
\renewcommand\linenumberfont{\normalfont\scriptsize\sffamily\color{black!55}}
\linenumbers
% header and footer
\usepackage{lastpage}
\usepackage{fancyhdr}
\pagestyle{fancy}
\fancyhf{}
\fancyhead[L]{\scriptsize\sffamily Review copy \textperiodcentered{} draft v5 \textperiodcentered{} printed 23 October 2026}
\fancyhead[R]{\scriptsize\sffamily page \thepage\ of \pageref{LastPage}}
\renewcommand{\headrulewidth}{0.3pt}
\fancypagestyle{plain}{\fancyhf{}\fancyhead[L]{\scriptsize\sffamily Review copy \textperiodcentered{} draft v5 \textperiodcentered{} printed 23 October 2026}\fancyhead[R]{\scriptsize\sffamily page \thepage\ of \pageref{LastPage}}}
% ruled notes column on every page
\usepackage{eso-pic,tikz}
\AddToShipoutPictureBG{%
  \begin{tikzpicture}[remember picture,overlay]
    \draw[black!35,line width=0.4pt]
      ([xshift=-5.85cm,yshift=-2.6cm]current page.north east) -- ([xshift=-5.85cm,yshift=2.6cm]current page.south east);
    \node[anchor=south west,inner sep=0pt,font=\scriptsize\sffamily,text=black!50]
      at ([xshift=-5.6cm,yshift=-2.5cm]current page.north east) {Notes};
    \foreach \y in {3.4,4.2,...,26.6}
      \draw[black!22,line width=0.3pt]
        ([xshift=-5.6cm,yshift=-\y cm]current page.north east) -- ([xshift=-0.9cm,yshift=-\y cm]current page.north east);
  \end{tikzpicture}}
% a ruled page for general comments at the end
\AtEndDocument{%
  \clearpage\nolinenumbers
  \section*{General comments}
  \begin{tikzpicture}[remember picture,overlay]
    \foreach \y in {5.0,5.8,...,26.6}
      \draw[black!22,line width=0.3pt]
        ([xshift=2.0cm,yshift=-\y cm]current page.north west) -- ([xshift=-6.2cm,yshift=-\y cm]current page.north east);
  \end{tikzpicture}}
% narrower column: let paragraphs stretch rather than overrun, and break three wide displays
\setlength{\emergencystretch}{3em}
\newcommand{\mkbreak}[1]{\\}
\newcommand{\mkalign}{\\}
\fi
\providecommand{\mkbreak}[1]{#1}% this copy: keeps the given spacing; review copy: line break
\providecommand{\mkalign}{&}%     this copy: next alignment column; review copy: new row

\title{Neighbourhood counting and new bounds for\\ small three-colour Ramsey numbers}
\author{Lachlan Dauth}
\date{23 October 2026}

\begin{document}
\maketitle

\begin{abstract}
We prove new upper and lower bounds for small three-colour Ramsey numbers.
Each upper bound $\R(t_1,t_2,t_3)\le n$ is proved by contradiction from a colouring of $K_n$ with no colour-$i$ $K_{t_i}$ for any $i$.
Ramsey bounds inside colour neighbourhoods restrict the degree profile of each vertex, its triple of colour degrees, to a finite list of admissible profiles.
A counting identity, or a flag-algebra certificate whose constraints include these local bounds, then shows that an explicit function of the profile, positive on every admissible profile, has nonpositive average, a contradiction.

By hand, merging the two triangle-free colours gives $\R(3,3,6)\le 88$, $\R(3,3,7)\le 129$ and $\R(3,3,8)\le 179$.
Also by hand, weighted triangle identities give $\R(3,4,5)\le 155$, $\R(3,4,6)\le 276$ and $\R(3,5,5)\le 346$, and four-vertex identities give $\R(4,4,4)\le 226$.
Five-vertex certificates with local caps, checked in exact arithmetic by computer, give $\R(4,4,4)\le 198$, $\R(3,4,4)\le 73$, $\R(3,4,5)\le 137$, $\R(3,5,5)\le 316$, $\R(3,4,6)\le 247$, $\R(3,4,7)\le 409$, $\R(3,5,6)\le 622$, $\R(3,4,8)\le 626$, $\R(3,3,6)\le 86$, $\R(3,3,7)\le 125$ and $\R(3,3,8)\le 174$.
Fed, together with four-vertex certificates, into the edge counting behind the hand bounds for $\R(3,3,k)$, they give $\R(3,3,8)\le 173$ and bounds for $\R(3,3,k)$ with $9\le k\le 15$, from $241$ to $892$.

The best previous upper bounds we know of are $229$, $77$, $157$ and $91$ for $\R(4,4,4)$, $\R(3,4,4)$, $\R(3,4,5)$ and $\R(3,3,6)$.
Radziszowski's survey lists no upper bound for the other numbers here.
Explicit colourings give $\R(3,3,k)\ge 133$, $183$, $203$, $232$, $267$ and $278$ for $k=9$, $11$, $12$, $13$, $14$ and $15$.
\end{abstract}

\begin{quote}\small
\textit{2020 Mathematics Subject Classification.} Primary 05C55; Secondary 05D10, 05C35, 90C22, 68V05.

\textit{Keywords.} Multicolour Ramsey numbers; neighbourhood recurrence; degree profiles; flag algebras; semidefinite programming; computer-assisted proof.
\end{quote}

\section{Introduction}

For integers $t_1,t_2,t_3\ge 1$, the Ramsey number $\R(t_1,t_2,t_3)$ is the least $n$ such that every colouring of the edges of $K_n$ with colours $1,2,3$ contains, for some $i$, a $K_{t_i}$ all of whose edges have colour $i$.
Two-colour numbers $\R(s,t)$ are defined in the same way.
We call colours $1$, $2$ and $3$ red, blue and green, respectively.

Renaming the colours shows that $\R(t_1,t_2,t_3)$ is symmetric in its arguments.
Up to this symmetry, only two values with every $t_i\ge 3$ are known exactly: $\R(3,3,3)=17$ \cite{GG55} and $\R(3,3,4)=30$ \cite{CFIM16}.
See Radziszowski's dynamic survey \cite{Rad26} for the current state of knowledge.

\paragraph{The neighbourhood recurrence.}
Most upper bounds for three-colour numbers come from the neighbourhood recurrence, implicit in \cite{GG55},
\begin{equation}\label{eq:recurrence}
\R(\bt)\;\le\;\R(\bt-\e_1)+\R(\bt-\e_2)+\R(\bt-\e_3)-1,
\end{equation}
where $\bt=(t_1,t_2,t_3)$ and $\e_i$ is the $i$th unit vector.

Write $N_i(v)$ for the colour-$i$ neighbourhood of a vertex $v$, $d_i(v)=|N_i(v)|$, and $\mathbf d(v)=(d_1(v),d_2(v),d_3(v))$ for its \emph{degree profile}.

In a colouring of $K_n$ with no colour-$i$ $K_{t_i}$ for any $i$, the set $N_i(v)$ has no colour-$i$ $K_{t_i-1}$, which would extend by $v$, and no colour-$j$ $K_{t_j}$ for $j\ne i$.
So
\[
d_i(v)\le\R(\bt-\e_i)-1,
\]
and summing over $i$ gives
\[
n-1\le\sum_i\big(\R(\bt-\e_i)-1\big),
\]
that is, $n$ is less than the right-hand side of \eqref{eq:recurrence}.

An entry equal to $2$ can be dropped, leaving a two-colour Ramsey number, as a colouring with no colour-$i$ $K_2$ has no colour-$i$ edges.
For instance \eqref{eq:recurrence} gives
\[
\R(3,3,k)\le 2\R(3,k)+\R(3,3,k-1)-1.
\]

\paragraph{Parity and other improvements.}
The recurrence has a \emph{parity refinement}: it is strict when its right-hand side is even and at least one term $\R(\bt-\e_i)$ is even \cite[6.1(a)]{Rad26}.
Indeed, at the order $n$ one less than the right-hand side, $n-1=\sum_i(\R(\bt-\e_i)-1)$, so equality holds in every cap $d_i(v)\le\R(\bt-\e_i)-1$.
For an $i$ with $\R(\bt-\e_i)$ even, every vertex then has odd colour-$i$ degree, which contradicts the handshake lemma as $n$ is odd.

Both statements, and every cap below, hold with the same proofs when a Ramsey number is replaced by a proven upper bound for it.
For example, $\R(3,3,6)\le 91$ \cite{Rad26,Boz26} follows from $\R(3,6)=18$ and $\R(3,3,5)\le 57$ (itself \eqref{eq:recurrence} with $\R(3,5)=14$ and $\R(3,3,4)=30$).
Here the right-hand side is $18+18+57-1=92$, and $92$ and $18$ are even.

Only a few three-colour upper bounds improve on \eqref{eq:recurrence} with parity \cite{Rad26}.
Among them are Boza's recent $\R(3,4,5)\le 157$ and $\R(4,4,4)\le 229$ \cite{Boz26}, obtained by an analogue of the parity argument modulo~$3$.
From $\R(4,5)=25$, $\R(3,3,5)\le 57$ and $\R(3,4,4)\le 77$, the recurrence gives $25+57+77-1=158$ and $3\cdot 77-1=230$ for these two numbers.
Parity does not apply, since every term is odd.
The survey also lists $\R(3,3,6)\le 91$, which \cite{Boz26} derives in the same way, among the bounds that improve on the recurrence with parity.
But, as shown above, it already follows from \eqref{eq:recurrence} with parity.
The bound $\R(3,4,4)\le 77$ is itself \eqref{eq:recurrence}, $18+30+30-1$, and has not been improved before.

\paragraph{Our approach.}
The recurrence bounds the colour neighbourhoods of one vertex at a time.
The parity refinement looks at all vertices together, but only at the order one less than the right-hand side of \eqref{eq:recurrence}, where every vertex has the \emph{maximal profile}
\[
(\R(\bt-\e_1)-1,\R(\bt-\e_2)-1,\R(\bt-\e_3)-1),
\]
and only at the parity of its entries.
Our arguments also look at all vertices together, but at smaller orders, where several profiles are possible, and they count small substructures instead.

\subsection{Results}

Our upper bounds for three-colour numbers are in Table~\ref{tab:upper}, and our lower bounds, with those listed in the survey \cite{Rad26}, in Table~\ref{tab:lower}.

Besides the exact values of $\R(3,3,3)$ and $\R(3,3,4)$, the survey lists upper bounds for only five three-colour numbers with every $t_i\ge 3$: $\R(3,3,5)$, $\R(3,3,6)$, $\R(3,4,4)$, $\R(3,4,5)$ and $\R(4,4,4)$.
We improve four of them, the bounds $\R(3,3,6)\le 91$, $\R(3,4,5)\le 157$ and $\R(4,4,4)\le 229$ recalled above and $\R(3,4,4)\le 77$.
We do not improve $\R(3,3,5)\le 57$.

\paragraph{Benchmarks.}
For the other numbers in Table~\ref{tab:upper} no upper bound is listed.
So the table gives instead a \emph{benchmark}, marked with a dagger: the result of applying \eqref{eq:recurrence} with parity repeatedly, starting from the values listed in \cite{Rad26,Boz26} and never using our new bounds.
For instance, $\R(3,7)=23$ gives $23+23+91-1=136$ for $\R(3,3,7)$, and then $\R(3,8)=28$ gives $28+28+136-1=191$ for $\R(3,3,8)$.
Parity improves neither.
We do not claim that any benchmark has appeared in the literature.

\paragraph{Hand and computer bounds.}
The bounds marked ``hand'' come from Sections~\ref{sec:three}, \ref{sec:types} and~\ref{sec:444}: double counting, which a reader can check with pen and paper.
Those marked ``computer'' come from Section~\ref{sec:five}: five-vertex certificates and, for $\R(3,3,k)$ with $k\ge 8$, the procedure of Section~\ref{sec:ladder} fed with exclusions and edge-count lower bounds from four- and five-vertex certificates.
The certificates are too large to check by hand but are verified in exact rational arithmetic by programs that share no code with the search that found them.

\paragraph{Companion paper.}
In a companion paper \cite{DC26} we derive from $\R(3,4,4)\le 73$ and $\R(4,4,4)\le 198$ new upper bounds for four-colour numbers, and give new bounds for two-colour numbers.


\begin{table}[t]
\centering\small
\begin{tabular}{lcccl}
\toprule
number & comparison & hand & computer & where\\
\midrule
$\R(3,3,6)$ & $91$ \cite{Rad26,Boz26} & $88$ & $86$ & Theorems~\ref{thm:336}, \ref{thm:flag678}\\
$\R(3,3,7)$ & $136^\dagger$ & $129$ & $125$ & Theorems~\ref{thm:337}, \ref{thm:flag678}\\
$\R(3,3,8)$ & $191^\dagger$ & $179$ & $173$ & Theorems~\ref{thm:338}, \ref{thm:flagladder}\\
$\R(3,3,k)$, $9\le k\le 15$ & $261,\dots,1019^\dagger$ & & Table~\ref{tab:flagladder} & Theorems~\ref{thm:ladder}, \ref{thm:flagladder}\\
$\R(3,4,4)$ & $77$ \cite{Rad26} & & $73$ & Theorem~\ref{thm:r34five}\\
$\R(3,4,5)$ & $157$ \cite{Boz26} & $155$ & $137$ & Theorems~\ref{thm:345}, \ref{thm:r34five}\\
$\R(3,4,6)$ & $287^\dagger$ & $276$ & $247$; $245^\ddagger$ & Theorems~\ref{thm:346}, \ref{thm:r34five}\\
$\R(3,5,5)$ & $359^\dagger$ & $346$ & $316$ & Theorems~\ref{thm:355}, \ref{thm:r34five}\\
$\R(3,4,7)$ & $479^\dagger$ & $456^\ddagger$ & $409$; $406^\ddagger$ & Proposition~\ref{prop:supp}, Theorem~\ref{thm:r34five}\\
$\R(3,4,8)$ & $748^\dagger$ & $700^\ddagger$ & $626$; $625^\ddagger$ & Proposition~\ref{prop:supp}, Theorem~\ref{thm:r34five}\\
$\R(3,5,6)$ & $730^\dagger$ & $683^\ddagger$ & $622$; $619^\ddagger$ & Proposition~\ref{prop:supp}, Theorem~\ref{thm:r34five}\\
$\R(4,4,4)$ & $229$ \cite{Boz26} & $226$ & $198$ & Theorems~\ref{thm:444a}, \ref{thm:r444five}\\
\bottomrule
\end{tabular}
\caption{Upper bounds for three-colour numbers.
$^\dagger$\,Benchmark where \cite{Rad26} lists no bound: \eqref{eq:recurrence} with parity, iterated from the values in \cite{Rad26,Boz26}, including the two-colour bounds marked $^*$ in Table~\ref{tab:inputs}, which the survey reports as unpublished.
Using these can only lower a benchmark.
For $9\le k\le 15$, the benchmarks are the row ``recurrence from $91$'' of Table~\ref{tab:ladder}.
The first row of that table gives the bounds of Theorem~\ref{thm:ladder}, which is computer-assisted but uses none of the flag-algebra certificates of Section~\ref{sec:five}, and Table~\ref{tab:flagladder} gives the bounds with the certificates of Section~\ref{sec:five}.
$^\ddagger$\,Uses starred bounds; in a cell ``$a$; $b^\ddagger$'', $a$ uses published inputs only.
With published inputs only, the hand arguments give $458$, $705$ and $685$, and the benchmarks in these rows become $483$, $757$ and $733$.
Of the other benchmarks, only $287$ uses a starred bound, $\R(4,6)\le 40$.
With the published $\R(4,6)\le 41$ it would be $288$.
Theorem~\ref{thm:444b} also gives $\R(4,4,4)\le 225$ from a four-vertex certificate with a single square.}
\label{tab:upper}
\end{table}

\begin{table}[t]
\centering
\begin{tabular}{lccl}
\toprule
number & lower bound in \cite{Rad26} & new lower bound & where\\
\midrule
$\R(3,3,9)$ & $129$ & $133$ & Theorem~\ref{thm:lower}\\
$\R(3,3,11)$ & $174$ & $183$ & Theorem~\ref{thm:lower}\\
$\R(3,3,12)$ & $194$ & $203$ & Theorem~\ref{thm:lower}\\
$\R(3,3,13)$ & $217$ & $232$ & Theorem~\ref{thm:lower}\\
$\R(3,3,14)$ & $242$ & $267$ & Theorem~\ref{thm:lower}\\
$\R(3,3,15)$ & $269$ & $278$ & Theorem~\ref{thm:ext15}\\
\bottomrule
\end{tabular}
\caption{Lower bounds for $\R(3,3,k)$ from explicit colourings.
For $k\le 14$, a known circulant colouring of $K_n$ showing $\R(5,k)\ge n+1$ has its $K_5$-free colour split into two triangle-free colours (Section~\ref{sec:lower}).
For $k=15$ eleven vertices are added to another split of the circulant colouring used for $k=14$ (Section~\ref{sec:extend}).
We have no new lower bound for other $\R(3,3,k)$.}
\label{tab:lower}
\end{table}

\subsection{Method}

Each of our upper bounds $\R(\bt)\le n$ is proved by contradiction from a colouring of $K_n$ with no colour-$i$ $K_{t_i}$ for any $i$, in three steps.
\begin{enumerate}[label=(\roman*)]
\item \emph{Local caps.}
By the proof of \eqref{eq:recurrence}, every profile is \emph{admissible}, that is, a triple of nonnegative integers with sum $n-1$ and $i$th entry at most $\R(\bt-\e_i)-1$.
If $n$ is $\delta$ less than the right-hand side of \eqref{eq:recurrence}, the shortfalls $\R(\bt-\e_i)-1-d_i(v)$, which are nonnegative, therefore sum to $\delta-1$.
So for small $\delta$ few profiles are admissible.
We also use \emph{mixed} caps on intersections of neighbourhoods of several vertices (Lemmas~\ref{lem:caps} and~\ref{lem:multicaps}), such as
\[
|N_i(v)\cap N_j(u)|\le\R(\bt-\e_i-\e_j)-1\qquad\text{for } u\in N_i(v).
\]
\item \emph{A counting identity, or a certificate.}
A small substructure is counted in two ways: triangles, two-edge paths or triples in Sections~\ref{sec:merge}--\ref{sec:ladder}; the nine triangle types that remain when red is triangle-free, in Section~\ref{sec:types}; four-vertex patterns, as probabilities for a random ordered quadruple of vertices, in Section~\ref{sec:444}.
One side depends only on the degree profiles.
The other involves terms such as the number of colour-$j$ edges inside $N_i(v)$, which the mixed caps bound from above (Corollary~\ref{cor:handshake}).
We merge two colours, or weight the triangle types, so that every term not determined by the profiles cancels or can be bounded in the required direction.
In Section~\ref{sec:five} the identity is replaced by a flag-algebra certificate \cite{Raz07}, mostly on five vertices: a nonnegative combination of caps, lower bounds on colour degrees and squares that leaves only terms determined by the profiles.
Either way we obtain
\[
\sum_v\Phi(\mathbf d(v))\le 0
\]
for an explicit function $\Phi$ of the profile.
\item \emph{A finite check.}
$\Phi$ is positive at every admissible profile, a contradiction.
\end{enumerate}

\paragraph{Relation to earlier work.}
The prototype is Goodman's count of monochromatic triangles in two-colourings \cite{Goo59}, based on the identity that the number of non-monochromatic triangles is $\frac12\sum_v d_1(v)d_2(v)$.
See \cite{Goo85} for three colours.

Bounding the triangles in the union of the triangle-free colours through common neighbourhoods is a standard step in asymptotic work on $\R(3,\dots,3,k)$ \cite{AR05,HW20,ACM26}.
Recently, \cite{ACM26} determined the order of magnitude of $\R(3,3,k)$.
What is new here is turning such counts into explicit bounds for small numbers.

The flag-algebra tables of Lidick\'y and Pfender \cite{LP21}, which use no local caps, give only $85$ for $\R(3,4,4)$ and $406$ for $\R(3,4,5)$, above the bounds from the recurrence.
Here the caps make the difference.

\paragraph{Using earlier bounds.}
The bounds build on each other: a new bound caps colour degrees in later proofs.
In Sections~\ref{sec:three} and~\ref{sec:ladder}, for $\R(3,3,k)$ with $k\ge 7$, the argument for $k-1$ enters the proof for $k$ twice:
\begin{itemize}
\item Its bound caps $d_3(v)$, for instance $d_3(v)\le 87$ in the proof of $\R(3,3,7)\le 129$.
\item Applied to the colouring on $N_3(v)$, which has no red or blue triangle and no green $K_{k-1}$, it bounds from below the number of red or blue edges inside $N_3(v)$ (Lemma~\ref{lem:line}(a)).
The density inequality (Lemma~\ref{lem:density}) counts vertex triples, and this number enters it on the side opposite the terms that the mixed caps bound from above, so it needs a lower bound.
\end{itemize}
The proof of Theorem~\ref{thm:336} needs no bound of this kind.

\subsection{What is proved by hand and what by computer}

\emph{Hand proofs.}
The bounds marked ``hand'' in Table~\ref{tab:upper} have complete proofs that can be checked by hand, once their inputs are granted.
Some inputs, such as $\R(3,3,4)=30$ \cite{CFIM16}, were themselves proved by computer.
Besides earlier results of this paper, these inputs are:
\begin{itemize}
\item the unstarred entries of Table~\ref{tab:inputs} (Section~\ref{sec:prelim});
\item the published two-colour bounds listed in its caption;
\item for the first three bounds of Proposition~\ref{prop:supp} only, the starred entries.
\end{itemize}

\emph{Computer-assisted proofs.}
The other upper bounds are computer-assisted: Theorems~\ref{thm:ladder} and~\ref{thm:444b}, and the results of Section~\ref{sec:five}.
Floating-point solvers were used only to find certificates and in the negative experiments we report.
Every certificate is verified in exact rational arithmetic.
Theorem~\ref{thm:444b} uses the computer only for an exact check of one linear identity among the probabilities of colour patterns on a random ordered quadruple of vertices.

\emph{Lower bounds.}
The lower bounds rest on explicit witnesses, checked by exact searches, and for $\R(3,3,15)$ also through Lemma~\ref{lem:cert}.

No upper-bound argument of ours, inputs aside, uses a SAT unsatisfiability result or a catalogue of critical colourings.
The checking programs are in the ancillary files, and Section~\ref{sec:verify} lists the few that use libraries beyond the standard ones.

\section{Preliminaries}\label{sec:prelim}

\paragraph{Notation.}
Let $\bt=(t_1,t_2,t_3)$.
A colouring of $E(K_n)$ is \emph{$\bt$-avoiding}, or a \emph{$(\bt;n)$-colouring} when we record the order, if for no $i$ is there a $K_{t_i}$ in colour $i$.
Restriction to fewer vertices preserves this property.
So a $(\bt;n)$-colouring exists if and only if $n<\R(\bt)$, and a vertex set on which a colouring is $\bt$-avoiding has at most $\R(\bt)-1$ vertices.

Besides $\R(2,a,b)=\R(a,b)$ from the introduction, so that $\R(2,2,a)=a$, we use $\R(1,a,b)=1$, as a single vertex is vacuously a $K_1$ in colour~$1$.
Such degenerate values arise in Lemmas~\ref{lem:caps} and~\ref{lem:multicaps} below when an entry drops to $2$ or $1$.

For a vertex set $S$, $e_j(S)$ is the number of colour-$j$ edges inside $S$.
It is unrelated to the unit vectors $\e_i$.

\paragraph{Proven bounds in place of Ramsey numbers.}
Our upper-bound arguments use a Ramsey number $\R(\mathbf s)$ only through the cap
\[
|S|\le\R(\mathbf s)-1
\]
for a set $S$ on which the colouring is $\mathbf s$-avoiding.
So a proven bound $\R(\mathbf s)\le B$ may replace $\R(\mathbf s)$ there, as in $\R(3,3,5)-1\le 56$.
Unless a proof says otherwise, $B$ is
\begin{itemize}
\item an unstarred entry of Table~\ref{tab:inputs},
\item one of the published bounds listed in its caption, or
\item a bound proved earlier in this paper.
\end{itemize}
The symbol $\R(\cdot)$ itself always denotes the true value.

\begin{lemma}[Local caps]\label{lem:caps}
Let $t_1,t_2,t_3\ge 3$, let $c$ be a $\bt$-avoiding colouring, $v$ a vertex and $i$ a colour.
Then $d_i(v)\le D_i:=\R(\bt-\e_i)-1$, and for every colour $j$ and every $u\in N_i(v)$,
\[
|N_i(v)\cap N_j(u)|\le C_{ij}:=\R(\bt-\e_i-\e_j)-1 .
\]
\end{lemma}

For instance, for $\bt=(3,3,k)$,
\begin{gather*}
D_1=D_2=\R(3,k)-1,\\
C_{12}=C_{21}=\R(2,2,k)-1=k-1\qquad\text{and}\qquad C_{11}=C_{22}=\R(1,3,k)-1=0,
\end{gather*}
the last restating that red and blue are triangle-free.

Where a number is given below for a cap $D_i$ or $C_{ij}$ whose Ramsey number is unknown, such as $D_3=56$ for $\bt=(3,3,6)$, it is the value $B-1$ from a bound $\R(\mathbf s)\le B$ as above.
The lemma holds with it since $\R(\mathbf s)\le B$.

\begin{proof}
Let $\mathbf s\le\bt$ entrywise, and let $S$ be a vertex set that contains, for each $l$ with $s_l<t_l$, no colour-$l$ $K_{s_l}$.
As $c$ is $\bt$-avoiding, its restriction to $S$ is then $\mathbf s$-avoiding, and so $|S|\le\R(\mathbf s)-1$.

For $S=N_i(v)$ and $\mathbf s=\bt-\e_i$, a colour-$i$ $K_{t_i-1}$ in $S$ would extend by $v$, so $d_i(v)\le D_i$.

For $S=N_i(v)\cap N_j(u)$ and $\mathbf s=\bt-\e_i-\e_j$, whose entries are positive as every $t_l\ge 3$:
\begin{itemize}
\item if $i\ne j$, a colour-$i$ $K_{t_i-1}$ in $S$ would extend by $v$, and a colour-$j$ $K_{t_j-1}$ by $u$;
\item if $i=j$, then $uv$ has colour $i$, so a colour-$i$ $K_{t_i-2}$ in $S$ would extend by both $u$ and $v$.
\end{itemize}
So $|S|\le C_{ij}$.
\end{proof}

\begin{corollary}\label{cor:handshake}
With $c$, $v$ and $i$ as in Lemma~\ref{lem:caps} and $j$ any colour,
\[
e_j(N_i(v))\le\lfloor C_{ij}\,d_i(v)/2\rfloor.
\]
\end{corollary}

\begin{proof}
Each $u\in N_i(v)$ has $|N_i(v)\cap N_j(u)|\le C_{ij}$ colour-$j$ neighbours in $N_i(v)$.
Apply the handshake lemma.
\end{proof}

The certificates of Section~\ref{sec:five} use caps on intersections of more than two neighbourhoods.
They come from the same argument.

\begin{lemma}[Caps on several neighbourhoods]\label{lem:multicaps}
Let $c$ be a $\bt$-avoiding colouring, $Y$ a nonempty set of vertices and $a$ a map from $Y$ to the colours.
For each colour $i$, let $w_i$ be the largest order of a colour-$i$ clique contained in $Y_i=\{y\in Y:a(y)=i\}$ (so $w_i=0$ if $Y_i$ is empty), and let $\mathbf w=(w_1,w_2,w_3)$.
Then $S=\bigcap_{y\in Y}N_{a(y)}(y)$ satisfies
\[
|S|\le\R(\bt-\mathbf w)-1.
\]
\end{lemma}

\begin{proof}
As $c$ is $\bt$-avoiding, $w_i\le t_i-1$, so every entry of $\bt-\mathbf w$ is at least $1$.
Every vertex of $S$ is joined in colour $i$ to every vertex of $Y_i$, and $S$ is disjoint from $Y$, as $y\notin N_{a(y)}(y)$.
So a colour-$i$ clique in $S$, together with a colour-$i$ clique of order $w_i$ in $Y_i$, is a colour-$i$ clique.
Hence $S$ contains no colour-$i$ $K_{t_i-w_i}$, the colouring on $S$ is $(\bt-\mathbf w)$-avoiding, and $|S|\le\R(\bt-\mathbf w)-1$.
\end{proof}

Lemma~\ref{lem:caps} is the case $|Y|\le 2$:
\begin{itemize}
\item $Y=\{v\}$ with $a(v)=i$ gives $\mathbf w=\e_i$ and the cap $D_i$; and
\item for $Y=\{v,u\}$ with $u\in N_i(v)$, $a(v)=i$ and $a(u)=j$, the vector $\mathbf w$ is $\e_i+\e_j$ if $j\ne i$, and $2\e_i$ if $j=i$, as $uv$ then has colour $i$.
\end{itemize}
When $w_i=t_i-1$ for some $i$, the cap is $\R(\bt-\mathbf w)-1=0$.

\paragraph{Triangle types and corners.}
In a $3$-colouring of $K_n$, write
\begin{itemize}
\item $M_i$ for the number of triangles with all edges of colour $i$,
\item $T_{i,j}$ ($i\ne j$) for the number with two edges of colour $i$, which meet at the \emph{apex} of the triangle, and one of colour $j$, and
\item $X$ for the number of rainbow triangles, whose edges have three different colours.
\end{itemize}
These count the ten \emph{types} of triangle.

A \emph{corner} at $v$ is a pair of edges $vu,vw$, with third edge $uw$.
It is an \emph{$\{i,j\}$-corner} if $vu$ and $vw$ have colours $i$ and $j$ in some order.
The number of $\{i,i\}$-corners at $v$ is $\binom{d_i(v)}{2}$, and that of $\{i,j\}$-corners with $j\ne i$ is $d_i(v)d_j(v)$.
For every colour $j$, including $j=i$, the $\{i,i\}$-corners at $v$ whose third edge has colour $j$ correspond to the colour-$j$ edges inside $N_i(v)$, so there are $e_j(N_i(v))$ of them.

On the other hand,
\begin{itemize}
\item a triangle of colour $i$ contains three $\{i,i\}$-corners;
\item one of type $T_{i,j}$ contains one $\{i,i\}$-corner, at its apex, whose third edge has colour $j$, and two $\{i,j\}$-corners; and
\item a rainbow triangle contains one $\{i,j\}$-corner for each of the three pairs $\{i,j\}$ with $i\ne j$.
\end{itemize}
Counting corners at the vertices and in the triangles gives the following identities.

\begin{lemma}[Corner identities]\label{lem:corners}
In every $3$-colouring of $K_n$, for every colour $i$,
\[
\sum_v\binom{d_i(v)}{2}=3M_i+\sum_{l\ne i}T_{i,l},\qquad 3M_i=\sum_v e_i(N_i(v)),
\]
and for every colour $j\ne i$,
\[
\sum_v d_i(v)d_j(v)=2T_{i,j}+2T_{j,i}+X,\qquad T_{i,j}=\sum_v e_j(N_i(v)).
\]
\end{lemma}

\begin{table}[htbp]
\centering\small
\begin{tabular}{lp{0.785\linewidth}}
\toprule
two colours & $\R(3,3)=6$, $\R(3,4)=9$, $\R(3,5)=14$, $\R(3,6)=18$, $\R(3,7)=23$, $\R(3,8)=28$, $\R(3,9)=36$, $\R(3,10)\le 41$ \cite{Ang25}, $\R(3,11)\le 50$ \cite{GR13}, $\R(3,12)\le 59$ \cite{Les}, $\R(3,13)\le 68$, $\R(3,14)\le 77$, $\R(3,15)\le 87$ \cite{GR13}; $\R(4,4)=18$, $\R(4,5)=25$, $\R(5,5)\le 46$ \cite{AM26}; and $\R(4,6)\le 40^*$, $\R(4,7)\le 58^*$, $\R(4,8)\le 79^*$, $\R(5,6)\le 85^*$\\
three colours & $\R(3,3,3)=17$ \cite{GG55}, $\R(3,3,4)=30$ \cite{CFIM16}; $\R(3,3,5)\le 57$ and $\R(3,4,4)\le 77$, both from \eqref{eq:recurrence}\\
\bottomrule
\end{tabular}
\caption{Inputs, as listed in \cite{Rad26}. The values $\R(3,k)$ are given for $3\le k\le 15$ since $\R(3,k)-1$ is the cap $D_1=D_2$ for $\bt=(3,3,k)$, and the procedure of Section~\ref{sec:ladder} runs through each of these $k$. $^*$\,Reported in \cite[Table~Ib]{Rad26} as unpublished (Angeltveit--McKay, personal communication); the published bounds are $\R(4,6)\le 41$, $\R(4,7)\le 61$, $\R(4,8)\le 84$ and $\R(5,6)\le 87$.}
\label{tab:inputs}
\end{table}

\section{Merging two triangle-free colours}\label{sec:merge}

\paragraph{Notation.}
In this section $k\ge 3$ and $n\ge 0$ are integers and $\bt=(3,3,k)$, so red and blue are triangle-free.
For a $(3,3,k;n)$-colouring:
\begin{itemize}
\item $H$ is its red-or-blue graph and $G=\overline H$ its green graph;
\item $h_v=d_1(v)+d_2(v)$ and $z_v=d_3(v)=n-1-h_v$ are the degrees of $v$ in $H$ and $G$;
\item $T(H)$ and $e(H)$ are the numbers of triangles and edges of $H$;
\item for a vertex set $S$, $e_H(S)=e_1(S)+e_2(S)$ is the number of edges of $H$ inside $S$.
\end{itemize}
By Lemma~\ref{lem:caps}, with $D:=D_1=D_2=\R(3,k)-1$ the common red and blue cap,
\[
d_1(v),d_2(v)\le D,\qquad\text{so}\qquad h_v\le 2D,\qquad\text{and}\qquad z_v\le D_3=\R(3,3,k-1)-1.
\]

Call an integer $h$ \emph{admissible} (as an $H$-degree) if
\[
\max(0,n-1-D_3)\le h\le\min(n-1,2D).
\]
Thus every $h_v$ is admissible.

A number $m_k(n)$ is a \emph{valid lower bound} if $e(H)\ge m_k(n)$ in every $(3,3,k;n)$-colouring.
We use this notion also for $k=2$, where $m_2(n)=\binom n2$ is valid, as a $(3,3,2;n)$-colouring has no green edge.

The lemmas below are also applied with $k-1$ and an order $z$ in place of $k$ and $n$, for instance to the colouring on a green neighbourhood.
Then $D$, $D_3$, admissibility and the functions below change accordingly.

\paragraph{Outline.}
We prove two inequalities,
\[
\sum_vF(h_v)\le 0\qquad\text{and}\qquad\sum_v\Psi(h_v)\le 0,
\]
in which the colouring enters only through the $H$-degrees (Lemmas~\ref{lem:commonnbr} and~\ref{lem:density}).
Both count paths with two edges in $H$ and bound $T(H)$ above by mixed caps (Lemma~\ref{lem:apex}).
\begin{itemize}
\item The first also uses the cap $C_{33}$ on common green neighbourhoods.
\item The second counts the vertex triples spanning one or two edges of $H$.
It needs instead a valid lower bound on the number of $H$-edges inside each green neighbourhood $N_3(v)$, which carries a $(3,3,k-1)$-colouring.
\end{itemize}
Lemma~\ref{lem:line} turns either inequality into a contradiction, excluding the order $n$, or into a valid lower bound $m_k(n)$, of the kind the second inequality needs at $k+1$.
Lemma~\ref{lem:delete} carries such bounds to larger orders.

\paragraph{The apex bound.}
In a $(3,3,k;n)$-colouring, every triangle of $H$ has two edges of one of red and blue and one of the other, so $T(H)=T_{1,2}+T_{2,1}$.
Counting each triangle at its apex by Lemma~\ref{lem:corners}, and applying Corollary~\ref{cor:handshake} with the mixed caps $C_{12}=C_{21}=\R(2,2,k)-1=k-1$, gives
\[
T(H)=\sum_v\big(e_2(N_1(v))+e_1(N_2(v))\big)\le\sum_v\Big(\Big\lfloor\tfrac{(k-1)d_1(v)}{2}\Big\rfloor+\Big\lfloor\tfrac{(k-1)d_2(v)}{2}\Big\rfloor\Big).
\]

\begin{lemma}[Apex bound]\label{lem:apex}
For $0\le h\le 2D$ let $A_k(h)=\lfloor(k-1)h/2\rfloor$, except that $A_k(2D)=(k-1)D-1$ when $k$ is even and $D$ is odd.
Then every $(3,3,k;n)$-colouring satisfies $T(H)\le\sum_vA_k(h_v)\le\frac{k-1}{2}\sum_vh_v$.
\end{lemma}

\begin{proof}
As $\lfloor x\rfloor+\lfloor y\rfloor\le\lfloor x+y\rfloor$, the summand at $v$ in the display above is at most $\lfloor(k-1)h_v/2\rfloor$.
If $h_v=2D$, the caps force $d_1(v)=d_2(v)=D$.
If also $k$ is even and $D$ is odd, then $(k-1)D$ is odd, and the summand is $2\lfloor(k-1)D/2\rfloor=(k-1)D-1$.
So the summand is at most $A_k(h_v)$, and the second inequality holds as $A_k(h)\le(k-1)h/2$.
\end{proof}

Write
\[
p_k(h)=3(k-1)h-6A_k(h),
\]
six times the gap between the two bounds of Lemma~\ref{lem:apex} at a vertex with $h_v=h$.
Thus $p_k=0$ for odd $k$.
For even $k$, $p_k(h)=3$ for odd $h$, $p_k(2D)=6$ if $D$ is odd, and $p_k(h)=0$ otherwise.

The counts below involve $6T(H)$.
Six times the first inequality of Lemma~\ref{lem:apex} reads
\begin{equation}\label{eq:apex6}
6T(H)\le\sum_v\big(3(k-1)h_v-p_k(h_v)\big).
\end{equation}
As $p_k\ge 0$, the term $p_k(h_v)$ may be dropped here, and hence the term $p_k(h)$, shown in brackets, from the functions of Lemmas~\ref{lem:commonnbr} and~\ref{lem:density}.
The hand proofs of Section~\ref{sec:three} drop it, and the procedure of Section~\ref{sec:ladder} (the \emph{ladder}) keeps it.

\begin{lemma}[Common-neighbour inequality]\label{lem:commonnbr}
Let $U\ge\R(3,3,k-2)$.
Every $(3,3,k;n)$-colouring satisfies $\sum_vF(h_v)\le 0$, where
\[
F(h)=F_{n,k,U}(h)=3h^2+(U+5-3n-3k)h+(n-U-1)(n-1)\;\big[\,+\,p_k(h)\,\big].
\]
With the bracket, $F$ depends also on $D$, through $p_k$.
\end{lemma}

\begin{proof}
\emph{Step 1: the cherry identity.}
A vertex $v$ is the middle vertex of $\binom{h_v}{2}$ \emph{cherries}, paths with two edges in $H$.
For vertices $x,y$ let $c_H(x,y)$ be the number of their common $H$-neighbours.
The endpoints of a cherry are joined either in $H$, closing a triangle of $H$, which contains three cherries, or by a green edge $xy$, and exactly $c_H(x,y)$ cherries have endpoints $x$ and $y$.
So
\[
\sum_v\binom{h_v}{2}=3T(H)+\sum_{xy\in E(G)}c_H(x,y).
\]

\emph{Step 2: common neighbours across a green edge.}
For a green edge $xy$, the $H$-neighbourhoods of $x$ and $y$ lie among the other $n-2$ vertices, and the vertices among these in neither form $N_3(x)\cap N_3(y)$.
By Lemma~\ref{lem:caps} with $i=j=3$,
\[
|N_3(x)\cap N_3(y)|\le C_{33}=\R(3,3,k-2)-1\le U-1,
\]
so inclusion--exclusion gives
\begin{equation}\label{eq:ie}
c_H(x,y)=h_x+h_y-(n-2)+|N_3(x)\cap N_3(y)|\le h_x+h_y+U+1-n.
\end{equation}

\emph{Step 3: summing over the green edges.}
Sum \eqref{eq:ie} over the green edges, using
\[
\sum_{xy\in E(G)}(h_x+h_y)=\sum_vh_vz_v\qquad\text{and}\qquad|E(G)|=\frac12\sum_vz_v,
\]
double the cherry identity, and bound $6T(H)$ by \eqref{eq:apex6}:
\[
\sum_vh_v(h_v-1)\le\sum_v\big(3(k-1)h_v-p_k(h_v)+2h_vz_v+(U+1-n)z_v\big).
\]

\emph{Step 4: collecting terms.}
Moving everything to the left and substituting $z_v=n-1-h_v$ turns the summand at $v$ into $F(h_v)$.
The coefficients of $h_v^2$ and $h_v$ are
\[
1+2=3\qquad\text{and}\qquad{-1}-3(k-1)-2(n-1)+(U+1-n)=U+5-3n-3k,
\]
and the constant term is $-(U+1-n)(n-1)$.
\end{proof}

\begin{lemma}[Density inequality]\label{lem:density}
For each admissible $h$, let $m_{k-1}(n-1-h)$ be a valid lower bound.
Then every $(3,3,k;n)$-colouring satisfies $\sum_v\Psi(h_v)\le 0$, where
\[
\Psi(h)=\Psi_{n,k}(h)=h\big(2h-n-3(k-1)\big)+2\,m_{k-1}(n-1-h)\;\big[\,+\,p_k(h)\,\big].
\]
Here $\Psi$ depends also on the chosen $m_{k-1}$ and, with the bracket, on $D$.
\end{lemma}

\begin{proof}
\emph{Step 1: two counting identities.}
Let $P_2$ and $P_1$ be the numbers of vertex triples spanning exactly two, respectively exactly one, edges of $H$.
A triple spanning two edges of $H$ contains exactly one cherry, whose endpoints are joined in green.
So the cherry identity in the proof of Lemma~\ref{lem:commonnbr} reads
\[
\sum_v\binom{h_v}2=3T(H)+P_2.
\]
There are $h_vz_v$ corners at $v$ with one edge in $H$ and one in $G$, and a triple contains two such corners if it spans one or two edges of $H$, and none otherwise.
So
\[
\sum_vh_vz_v=2P_1+2P_2.
\]
Twice the first identity minus the second, with $h_v-1-z_v=2h_v-n$, gives
\[
6T(H)=\sum_vh_v(2h_v-n)+2P_1 .
\]

\emph{Step 2: the triples with one edge of $H$.}
In a triple with exactly one edge of $H$, the vertex $v$ opposite that edge is joined to both its ends in green, so the edge lies inside $N_3(v)$.
Conversely, each edge of $H$ inside $N_3(v)$ forms such a triple with $v$.
So
\[
P_1=\sum_ve_H(N_3(v)).
\]
The colouring on $N_3(v)$ is a $(3,3,k-1;z_v)$-colouring, as in the proof of Lemma~\ref{lem:caps}, and $h_v$ is admissible, so $e_H(N_3(v))\ge m_{k-1}(z_v)$.

\emph{Step 3: comparison with the apex bound.}
Hence
\[
6T(H)\ge\sum_v\big(h_v(2h_v-n)+2m_{k-1}(z_v)\big),
\]
and comparing with \eqref{eq:apex6} gives $\sum_v\Psi(h_v)\le 0$.
\end{proof}

\begin{lemma}\label{lem:line}
Let $f$ be a function such that every $(3,3,k;n)$-colouring satisfies $\sum_vf(h_v)\le 0$.
\begin{enumerate}[label=(\alph*)]
\item If $\beta>0$ and $f(h)\ge\alpha-\beta h$ for every admissible $h$, then every $(3,3,k;n)$-colouring has $e(H)\ge n\alpha/(2\beta)$.
Hence $\lceil n\alpha/(2\beta)\rceil$ is a valid lower bound $m_k(n)$.
\item If $n\ge 1$ and $f(h)>0$ for every admissible $h$, then no $(3,3,k;n)$-colouring exists.
\end{enumerate}
\end{lemma}

\begin{proof}
Take a $(3,3,k;n)$-colouring.
Every $h_v$ is admissible, and $\sum_vh_v=2e(H)$.

(a) $0\ge\sum_vf(h_v)\ge\sum_v(\alpha-\beta h_v)=n\alpha-2\beta e(H)$.
Divide by $2\beta>0$, and round up since $e(H)$ is an integer.

(b) $\sum_vf(h_v)$ is a sum of $n\ge 1$ positive terms, so it is positive, a contradiction.
\end{proof}

Part~(a) bounds the average $H$-degree $2e(H)/n$ below by $\alpha/\beta$, the zero of $\alpha-\beta h$.
So we want a line below $f$ at the admissible integers whose zero lies as far right as possible.
Section~\ref{sec:three} uses chords (Lemmas~\ref{lem:m6more} and~\ref{lem:m7}), and Section~\ref{sec:ladder} edges of the lower convex hull of the points $(h,f(h))$.
Part~(b) is the finite check, step~(iii) of the method.

Lemma~\ref{lem:delete} passes valid lower bounds to larger orders,
\begin{itemize}
\item by deleting a vertex, whose $H$-degree is bounded below as its green degree is capped, or
\item by averaging over the vertex subsets of a fixed size.
\end{itemize}
So one application of Lemma~\ref{lem:line} can serve several orders.

\begin{lemma}\label{lem:delete}
Let $n\ge 1$, and for each $n'<n$ let $m_k(n')$ be a valid lower bound.
Then each of the following is also a valid lower bound for order $n$:
\begin{enumerate}[label=(\alph*)]
\item $m_k(n-1)+\max\{0,\,n-\R(3,3,k-1)\}$;
\item $\big\lceil m_k(s)\,n(n-1)/(s(s-1))\big\rceil$, for any $2\le s<n$.
\end{enumerate}
\end{lemma}

\begin{proof}
Take a $(3,3,k;n)$-colouring.

(a) By Lemma~\ref{lem:caps}, each vertex $v$ has green degree $n-1-h_v\le\R(3,3,k-1)-1$, so $h_v\ge\max\{0,\,n-\R(3,3,k-1)\}$.
Deleting any vertex $v$ leaves a $(3,3,k;n-1)$-colouring with red-or-blue graph $H-v$, so
\[
e(H)=e(H-v)+h_v\ge m_k(n-1)+\max\{0,\,n-\R(3,3,k-1)\}.
\]

(b) Each $s$-subset $S$ of the vertices induces a $(3,3,k;s)$-colouring, so $e(H[S])\ge m_k(s)$.
Summing over the $\binom ns$ subsets counts each edge of $H$ once for each of the $\binom{n-2}{s-2}$ subsets containing both its ends, so
\[
e(H)\ge m_k(s)\binom ns\Big/\binom{n-2}{s-2}=m_k(s)\,n(n-1)/(s(s-1)).
\]
Round up, as $e(H)$ is an integer.
\end{proof}

\section{Three upper bounds by hand}\label{sec:three}

Theorem~\ref{thm:336} uses Lemma~\ref{lem:commonnbr} alone.
Applying the same lemma to $(3,3,6;z)$-colourings, with chords in Lemma~\ref{lem:line}(a), gives valid lower bounds $m_6(z)$ for $79\le z\le 87$ (Lemma~\ref{lem:m6more}).
With these, Lemma~\ref{lem:density} gives Theorem~\ref{thm:337} and, with Lemma~\ref{lem:delete}(a), the bounds $m_7(z)$ for $124\le z\le 128$ of Lemma~\ref{lem:m7}, which give Theorem~\ref{thm:338} in turn.

By computer, Section~\ref{sec:five} lowers the three bounds to $86$, $125$ and $173$ (Theorems~\ref{thm:flag678} and~\ref{thm:flagladder}).
The certificates of Theorem~\ref{thm:r34five} use the bounds proved here as caps.

\begin{theorem}\label{thm:336}
$\R(3,3,6)\le 88$.
\end{theorem}

\begin{proof}
\emph{Step 1: the degree cap.}
Suppose a $(3,3,6;88)$-colouring exists.
Here $D=\R(3,6)-1=17$, so $h_v\le 2D=34$.

\emph{Step 2: the common-neighbour inequality.}
Lemma~\ref{lem:commonnbr} with $n=88$, $k=6$ and $U=\R(3,3,4)=30$, without the bracket, gives
\[
F(h)=3h^2-247h+4959.
\]

\emph{Step 3: positivity.}
Writing $x=34-h$, we get
\[
F=29+43x+3x^2,
\]
which is positive for every $x\ge 0$, that is, for every $h\le 34$, in particular at every admissible $h$ (the green cap is not needed).
Lemma~\ref{lem:line}(b) gives a contradiction.
\end{proof}

\paragraph{Why the argument stops here.}
The argument stops at $88$.
At $n=87$,
\[
F_{87,6,30}(h)=3h^2-244h+4816
\]
has $F_{87,6,30}(34)=-12$, or $-6$ with the bracket $p_6(34)=6$ (as $34=2D$ with $D$ odd).
And $h=34$ is still admissible, as the green degree $86-34=52$ is within the cap $D_3=56$.

Nor does the procedure of Section~\ref{sec:ladder}, which also applies the density inequality, exclude $87$ (Remark~\ref{rem:sanity}).
The certificates of Section~\ref{sec:five} do (Theorem~\ref{thm:flag678}).

\paragraph{Lower bounds $m_6(z)$ from chords.}
For $(3,3,6;z)$-colourings, Lemma~\ref{lem:commonnbr} with $n=z$, $k=6$ and $U=\R(3,3,4)=30$, without the bracket, gives $\sum_vF_{z,6,30}(h_v)\le 0$ with
\[
F_{z,6,30}(h)=3h^2-(3z-17)h+(z-31)(z-1).
\]
As Lemma~\ref{lem:line}(a) needs a line below $F_{z,6,30}$ only at integers, we use a chord rather than a tangent.
For an integer $q$,
\begin{equation}\label{eq:chord}
\begin{gathered}
F_{z,6,30}(h)=\alpha-\beta h+3(h-q)(h-q-1)\ \ge\ \alpha-\beta h,\\
\beta=3z-20-6q,\qquad \alpha=(z-31)(z-1)-3q(q+1),
\end{gathered}
\end{equation}
where the inequality holds at every integer $h$, as $(h-q)(h-q-1)$ is then a product of two consecutive integers and so nonnegative.
So if $\beta>0$, Lemma~\ref{lem:line}(a) shows that $\lceil z\alpha/(2\beta)\rceil$ is a valid lower bound $m_6(z)$.

The chord $\alpha-\beta h$ meets $F_{z,6,30}$ at $h=q$ and $h=q+1$.
At the smaller root $r$ of $F_{z,6,30}$ it takes the value $-3(r-q)(r-q-1)$, which is positive exactly when $q<r<q+1$.
So for $q=\lfloor r\rfloor$ the chord, if decreasing ($\beta>0$), crosses zero to the right of $r$ ($r$ is not an integer for $79\le z\le 87$, as the discriminant $(3z-17)^2-12(z-31)(z-1)$, which decreases from $3472$ to $1744$ on this range, is never a perfect square).
By contrast, for every other $q$ with $\beta>0$, and for every decreasing tangent, which is at most $F_{z,6,30}(r)=0$ at $r$, the zero lies at or left of $r$.
Thus $q=\lfloor r\rfloor$ gives the strongest of these bounds.
The procedure of Section~\ref{sec:ladder} raises some of them (Remark~\ref{rem:sanity}).

\begin{lemma}\label{lem:m6more}
For $79\le z\le 87$,
\[
m_6(z)=\lceil z\alpha/(2\beta)\rceil
\]
is a valid lower bound for $e(H)$ in $(3,3,6;z)$-colourings, with $\alpha$ and $\beta$ from \eqref{eq:chord} for $q=\lfloor r\rfloor$, where $r$ is the smaller root of $F_{z,6,30}$, shown rounded to two decimals:
\begin{center}\small\upshape
\begin{tabular}{c*{9}{c}}
\toprule
$z$ &79&80&81&82&83&84&85&86&87\\
$r$ &26.85&27.63&28.42&29.24&30.08&30.94&31.83&32.75&33.71\\
$q$ &26&27&28&29&30&30&31&32&33\\
$\beta$ &61&58&55&52&49&52&49&46&43\\
$\alpha$ &1638&1603&1564&1521&1474&1609&1560&1507&1450\\
$m_6(z)$ &1061&1106&1152&1200&1249&1300&1354&1409&1467\\
\bottomrule
\end{tabular}
\end{center}
\end{lemma}

\begin{proof}
Every $\beta$ in the table is positive, so \eqref{eq:chord} and Lemma~\ref{lem:line}(a) apply.
\end{proof}

\begin{theorem}\label{thm:337}
$\R(3,3,7)\le 129$.
\end{theorem}

\begin{proof}
\emph{Step 1: the admissible degrees.}
Suppose a $(3,3,7;129)$-colouring exists.
Here $h_v\le 2D=2(\R(3,7)-1)=44$, and $z_v=128-h_v\le 87$ by Theorem~\ref{thm:336}.
So the admissible $H$-degrees are $41,\dots,44$, with green degrees $84,\dots,87$.

\emph{Step 2: the density inequality.}
Lemma~\ref{lem:density} with $n=129$, $k=7$ and the values $m_6(z)$ of Lemma~\ref{lem:m6more} gives $\sum_v\Psi(h_v)\le 0$ for
\[
\Psi(h)=h(2h-147)+2m_6(128-h).
\]
The bracket $p_7$ vanishes, as $k$ is odd.

\emph{Step 3: positivity.}
But $\Psi$ is positive at every admissible $h$:
\begin{center}
\begin{tabular}{cccc}
\toprule
$h$ & $z=128-h$ & $m_6(z)$ & $\Psi(h)$\\
\midrule
41&87&1467&269\\ 42&86&1409&172\\ 43&85&1354&85\\ 44&84&1300&4\\
\bottomrule
\end{tabular}
\end{center}
so Lemma~\ref{lem:line}(b) gives a contradiction.
\end{proof}

\paragraph{Comparison.}
With Theorem~\ref{thm:336}, \eqref{eq:recurrence} gives only $23+23+88-1=133$ for $\R(3,3,7)$.

The common-neighbour inequality does not exclude $n=129$, and neither does the density inequality without a nontrivial lower bound $m_6$:
\begin{itemize}
\item Lemma~\ref{lem:commonnbr}, with $U=57$ from $\R(3,3,5)\le 57$, gives $F(h)=3h^2-346h+9088$ ($p_7=0$), which is decreasing for $h\le 57$, so its largest value at an admissible $h$ is $F(41)=-55$;
\item and Lemma~\ref{lem:density} with the trivial $m_6=0$ gives $\Psi(h)=h(2h-147)<0$ at every admissible $h$.
\end{itemize}

By Lemma~\ref{lem:m6more}, in contrast, each green neighbourhood, which carries a $(3,3,6;z)$-colouring with $84\le z\le 87$, has average red-or-blue degree at least $2m_6(z)/z>30$, close to the cap $2(\R(3,6)-1)=34$.

\begin{lemma}\label{lem:m7}
Valid lower bounds for $e(H)$ in $(3,3,7;z)$-colourings are
\begin{gather*}
m_7(124)=2495,\qquad m_7(125)=2532,\qquad m_7(126)=2570,\\
m_7(127)=2609\qquad\text{and}\qquad m_7(128)=2649.
\end{gather*}
\end{lemma}

\begin{proof}
\emph{Step 1: the admissible degrees.}
Take a $(3,3,7;124)$-colouring.
Here $h_v\le 2(\R(3,7)-1)=44$, and $z_v=123-h_v\le 87$ by Theorem~\ref{thm:336}.
So the admissible $H$-degrees are $36,\dots,44$, with green degrees $79,\dots,87$, those covered by Lemma~\ref{lem:m6more}.

\emph{Step 2: the density inequality.}
Lemma~\ref{lem:density} with $n=124$, $k=7$ and these values $m_6(z)$ gives $\sum_v\Psi(h_v)\le 0$ for
\[
\Psi(h)=h(2h-142)+2m_6(123-h).
\]
The bracket $p_7$ vanishes as $k$ is odd.

\emph{Step 3: a chord below $\Psi$.}
As $\Psi$ takes both signs, we apply Lemma~\ref{lem:line}(a) with the line $3138-78h$, the chord of $\Psi$ through $h=40$ and $h=41$, where $\Psi$ changes sign:
\begin{center}\small
\begin{tabular}{c*{9}{r}}
\toprule
$h$ &36&37&38&39&40&41&42&43&44\\
$\Psi(h)$ &414&302&200&104&18&$-60$&$-132$&$-196$&$-254$\\
$\Psi(h)-(3138-78h)$ &84&50&26&8&0&0&6&20&40\\
\bottomrule
\end{tabular}
\end{center}
The last row is nonnegative, so
\[
\Psi(h)\ge 3138-78h
\]
at every admissible $h$.
Lemma~\ref{lem:line}(a) with $\alpha=3138$ and $\beta=78$ then gives
\[
e(H)\ge 124\cdot 3138/156>2494.
\]
So $m_7(124)=2495$ is valid.

\emph{Step 4: the larger orders.}
For $125\le z\le 128$, Lemma~\ref{lem:delete}(a) with $k=7$ and $\R(3,3,6)\le 88$ gives
\[
m_7(z)=m_7(z-1)+(z-88),
\]
adding $37$, $38$, $39$ and $40$ in turn.
\end{proof}

\begin{theorem}\label{thm:338}
$\R(3,3,8)\le 179$.
\end{theorem}

\begin{proof}
\emph{Step 1: the admissible degrees.}
Suppose a $(3,3,8;179)$-colouring exists.
Here $h_v\le 2D=2(\R(3,8)-1)=54$, and $z_v=178-h_v\le 128$ by Theorem~\ref{thm:337}.
So the admissible $H$-degrees are $50,\dots,54$, with green degrees $124,\dots,128$, exactly the orders of Lemma~\ref{lem:m7}.

\emph{Step 2: the density inequality.}
Lemma~\ref{lem:density} with $n=179$, $k=8$ and these values $m_7(z)$, without the bracket, gives $\sum_v\Psi(h_v)\le 0$ for
\[
\Psi(h)=h(2h-200)+2m_7(178-h).
\]

\emph{Step 3: positivity.}
But $\Psi$ is positive at all five admissible $h$:
\begin{center}
\begin{tabular}{cccc}
\toprule
$h$ & $z=178-h$ & $m_7(z)$ & $\Psi(h)$\\
\midrule
50&128&2649&298\\ 51&127&2609&220\\ 52&126&2570&148\\ 53&125&2532&82\\ 54&124&2495&22\\
\bottomrule
\end{tabular}
\end{center}
So Lemma~\ref{lem:line}(b) gives a contradiction.
\end{proof}

\section{The ladder: \texorpdfstring{$\R(3,3,k)$}{R(3,3,k)} for \texorpdfstring{$9\le k\le 15$}{9 <= k <= 15}}\label{sec:ladder}

The procedure of this section, the \emph{ladder}, automates Section~\ref{sec:three}, in levels $k=3,4,\dots,15$, each level using the bounds and tables of the levels below it.
Write $B_k$ for the upper bound on $\R(3,3,k)$ that it uses.
Level $k$ records valid lower bounds $m_k(n)$ for $2\le n\le B_k-1$.
For $k\le 5$, $B_k$ is an input: $B_1=\R(3,3,1)=1$ and $B_2=\R(3,3,2)=6$ are degenerate values (Section~\ref{sec:prelim}), and $B_3=17$, $B_4=30$ and $B_5=57$, from Table~\ref{tab:inputs}, are the only three-colour inputs.
For $k\ge 6$, $B_k$ is the first order excluded at level $k$, where the level stops; so $\R(3,3,k)\le B_k$.

\paragraph{Each level.}
The procedure starts from the trivial table $m_2(s)=\binom s2$ for $s\le 5$ (Section~\ref{sec:merge}), and puts $m_k(0)=m_k(1)=0$.
Level $k$ uses:
\begin{itemize}
\item the caps $D=r_k-1$, where $r_k$ is the value or bound for $\R(3,k)$ in Table~\ref{tab:inputs} (so $D=40,49,58,67,76,86$ for $k=10,\dots,15$), and $D_3=B_{k-1}-1$;
\item the value $U=B_{k-2}$ in $F$ (Lemma~\ref{lem:commonnbr});
\item and in $\Psi$ (Lemma~\ref{lem:density}) the table $m_{k-1}$, which covers every green degree $z\le D_3$.
\end{itemize}
Both $F$ and $\Psi$ include the bracket $p_k$.

The level treats the orders $n=2,3,\dots$ from the smallest upwards, so that Lemma~\ref{lem:delete} can use the values of $m_k$ at smaller orders.
Levels $3$, $4$ and $5$ run through all $n<B_k$ and exclude none (Remark~\ref{rem:sanity}).
They serve to build $m_3$, $m_4$ and $m_5$, the last of which enters $\Psi$ at level~$6$.

\paragraph{Treatment of one order.}
At order $n$, the admissible $H$-degrees are the integers $h$ with
\[
h_{\min}\le h\le\min(2D,n-1),
\]
where $h_{\min}=\max(0,n-B_{k-1})$.
\begin{enumerate}[label=(\roman*)]
\item If $F$ or $\Psi$ is positive at every admissible $h$, then $n$ is excluded by Lemma~\ref{lem:line}(b).
\item Otherwise, for each $f\in\{F,\Psi\}$ let $h^*_f$ be the least real $h\ge h_{\min}$ at which the lower convex hull of the points $(h,f(h))$, $h$ admissible, is at most $0$.
Then $\lceil nh^*_f/2\rceil$ is a valid lower bound.
\begin{itemize}
\item If $h^*_f=h_{\min}$, that is, $f(h_{\min})\le 0$, this holds because every $H$-degree is at least $h_{\min}$.
\item Otherwise the hull edge reaching $(h^*_f,0)$ from the left lies on a line $\alpha-\beta h$ with $\beta>0$ and zero at $h^*_f$.
This line lies below the hull, so it is a \emph{supporting line}:
\[
f(h)\ge\alpha-\beta h
\]
at every admissible $h$.
Lemma~\ref{lem:line}(a) then gives $\lceil n\alpha/(2\beta)\rceil=\lceil nh^*_f/2\rceil$ (no supporting line with $\beta>0$ does better, as each lies below the hull and so is at most $0$ at $h^*_f$).
\end{itemize}
\item The recorded $m_k(n)$ is the largest of the bounds from $F$ and $\Psi$ and those that Lemma~\ref{lem:delete}(a) and~(b) give from the values at smaller orders, with $B_{k-1}$ in place of $\R(3,3,k-1)$ in~(a).
\end{enumerate}

\begin{remark}[Sanity checks]\label{rem:sanity}
The levels $k\le 8$ of the procedure behave as follows.

For $k=3$ and $k=4$ it excludes no order below the exact values $\R(3,3,3)=17$ and $\R(3,3,4)=30$.
As colourings exist at all these orders, this tests soundness.

At $n=16$ the caps $D=\R(3,3)-1=5$ and $D_3=\R(3,3,2)-1=5$ give $h_v\le 10$ and $h_v=15-z_v\ge 10$.
So $h_v=10$ at every vertex, and $e(H)=80$, in every $(3,3,3;16)$-colouring, such as that of Greenwood and Gleason \cite{GG55}, in which each colour class is the $5$-regular Clebsch graph.
Lemmas~\ref{lem:commonnbr} and~\ref{lem:density}, with $U=\R(3,3,1)=1$ and the valid $m_2(s)=\binom s2$ (a $(3,3,2)$-colouring has no green edge), give
\[
F_{16,3,1}(h)=3h^2-51h+210\qquad\text{and}\qquad\Psi_{16,3}(h)=h(2h-22)+2\binom{15-h}{2},
\]
with $p_3=0$, and both vanish at $h=10$.
So neither inequality excludes $n=16$, and the procedure records the exact value $m_3(16)=80$.
This tests the correctness of the inequalities rather than their strength.

For $k=5$ it excludes no order below $57$, the least order at which the caps alone leave no admissible $h$, so it does not improve $\R(3,3,5)\le 57$.

For $k=6,7,8$ it reaches exactly the hand-proved bounds $88$, $129$ and $179$, excluding no smaller order, although it uses more than those proofs do:
\begin{itemize}
\item the bracket $p_k$;
\item both inequalities at every order;
\item supporting lines from lower convex hulls;
\item and Lemma~\ref{lem:delete}(b).
\end{itemize}
These raise some values $m_k(z)$ above those of the hand proofs, but not the bounds.
\end{remark}

\begin{theorem}[computer-assisted]\label{thm:ladder}
The procedure gives
\[
(B_9,\dots,B_{15})=(247,315,407,512,628,756,901).
\]
Hence $\R(3,3,k)\le B_k$ for $9\le k\le 15$.
These bounds form the first row of Table~\ref{tab:ladder}.
\end{theorem}

The run recomputes $B_6=88$, $B_7=129$ and $B_8=179$, agreeing with Section~\ref{sec:three}, so the theorem does not rest on the hand proofs there.
Its external inputs, the values and bounds for $\R(3,k)$ and $B_3,B_4,B_5$ from Table~\ref{tab:inputs}, are all unstarred, so it uses no unpublished bound.

For $9\le k\le 15$, it is $\Psi$ that excludes $B_k$ by Lemma~\ref{lem:line}(b): its minimum over the admissible $h$ at $n=B_k$, in the last row of Table~\ref{tab:ladder}, is positive, whereas $F$ is negative at $h=2D$, the largest admissible $h$.
In every column this minimum is attained at $h=2D$, where $p_k(2D)=0$, so it equals
\[
2D\big(4D-B_k-3(k-1)\big)+2m_{k-1}(z),\qquad z=B_k-1-2D.
\]
For $k=9,\dots,15$ the green degree is $z=176$, $234$, $308$, $395$, $493$, $603$, $728$, and $m_{k-1}(z)=4599$, $7317$, $11810$, $18216$, $26598$, $37421$, $51609$.

The run writes a certificate, which a checker that shares no code with the generating program verifies in exact rational arithmetic.
With valid lower bounds and exclusions from the certificates of Section~\ref{sec:five}, the same procedure gives the smaller bounds of Table~\ref{tab:flagladder} (Theorem~\ref{thm:flagladder}).

\begin{table}[htbp]
\centering\small
\begin{tabular}{l*{7}{c}}
\toprule
$k$ & 9 & 10 & 11 & 12 & 13 & 14 & 15\\
\midrule
Theorem~\ref{thm:ladder} & 247 & 315 & 407 & 512 & 628 & 756 & 901\\
recurrence from $91^\dagger$ & 261 & 342 & 441 & 558 & 693 & 846 & 1019\\
recurrence from $B_{k-1}$ & 249 & 328 & 413 & 524 & 647 & 781 & 929\\
admissible $h$ & 68--70 & 68--80 & 92--98 & 105--116 & 116--134 & 128--152 & 145--172\\
minimum of $\Psi$ & 28 & 74 & 2 & 124 & 132 & 210 & 190\\
\bottomrule
\end{tabular}
\caption{The ladder.
The row ``recurrence from $91$'' gives the daggered benchmarks of Table~\ref{tab:upper}: \eqref{eq:recurrence} with parity, iterated from $\R(3,3,6)\le 91$ \cite{Rad26} through $136$ and $191$ for $k=7,8$.
The row ``recurrence from $B_{k-1}$'' applies \eqref{eq:recurrence} with parity once, to $B_{k-1}$.
The last two rows concern the excluded order $n=B_k$: its admissible $H$-degrees, and the minimum over them of $\Psi$, with the bracket $p_k$ and the table $m_{k-1}$ recorded at level $k-1$.
The tables $m_{k-1}$ are not reproduced.
The programs for Theorem~\ref{thm:ladder} (Section~\ref{sec:verify}) regenerate them.}
\label{tab:ladder}
\end{table}

\begin{proof}[Verification of Theorem~\ref{thm:ladder}]
\emph{The certificate.}
The certificate has one row for each level $k$ with $3\le k\le 15$ and each order $2\le n\le B_k-1$, so $\sum_{k=3}^{15}(B_k-2)=4240$ rows.
Each row states a claimed value of $m_k(n)$ and, for each $f\in\{F,\Psi\}$, certifies the bound from $f$:
\begin{itemize}
\item either it records the admissible ends $h'<h''$ of the hull edge used, so that the supporting line $\alpha-\beta h$ is the line through $(h',f(h'))$ and $(h'',f(h''))$;
\item or it records that $f(h_{\min})\le 0$, and the bound is then $\lceil nh_{\min}/2\rceil$, which needs no line, as every $H$-degree is at least $h_{\min}$.
\end{itemize}

\emph{The checker.}
A checker that does not import the generating program treats the rows in the order of the procedure, in exact rational arithmetic.
For each recorded line it evaluates $f$ from its definition, with the bracket $p_k$ and, in $\Psi$, the values $m_{k-1}$ it has already verified.
It then checks the hypotheses of Lemma~\ref{lem:line}(a): $\beta>0$, and $f(h)\ge\alpha-\beta h$ at every admissible $h$.
Summed over the rows, there are $420{,}545$ admissible degrees.
It accepts the claimed $m_k(n)$ only if it is at most the largest of the two certified bounds and those that Lemma~\ref{lem:delete}(a) and~(b) give from the values it has verified at smaller orders.

\emph{The bounds $B_k$.}
Finally, for $k\le 5$ it checks that $B_k$ is the input, and for $6\le k\le 15$ that at $n=B_k$, which has no row, one of $F$ and $\Psi$ is positive at every admissible $h$, so that Lemma~\ref{lem:line}(b) excludes $B_k$.
\end{proof}

\section{Weighted triangle identities}\label{sec:types}

In this section $t_1=3$ and $t_2,t_3\ge 4$, so red is the only colour required to be triangle-free, and $M_1=0$.
This covers $\R(3,4,5)$, $\R(3,4,6)$ and $\R(3,5,5)$, and the three numbers of Proposition~\ref{prop:supp}.
Theorem~\ref{thm:r34five} improves all six bounds of this section by computer.
The arguments here can be checked by hand.

The merging of Section~\ref{sec:merge} needs two triangle-free colours, so here we keep the three colours apart and take a weighted sum of the corner identities of Lemma~\ref{lem:corners}.

\begin{lemma}\label{lem:weighted}
For real weights $w_1,w_2,w_3,w_{12},w_{13},w_{23}$, not necessarily positive, put $w_{ji}:=w_{ij}$ and
\[
W(\mathbf d)=\sum_i w_i\binom{d_i}{2}+\sum_{i<j}w_{ij}d_id_j .
\]
Then for every $3$-colouring of $K_n$,
\[
\sum_vW(\mathbf d(v))=\sum_i 3w_iM_i+\sum_{i\ne j}(w_i+2w_{ij})T_{i,j}+(w_{12}+w_{13}+w_{23})X .
\]
\end{lemma}

\begin{proof}
In Lemma~\ref{lem:corners}, multiply the identity for $\sum_v\binom{d_i(v)}2$ by $w_i$ and that for $\sum_v d_i(v)d_j(v)$ by $w_{ij}$, and sum over all $i$ and all pairs $i<j$.
The first gives $3w_i$ to $M_i$ and $w_i$ to each $T_{i,l}$ with $l\ne i$.
The second, for the pair $i<j$, gives $2w_{ij}$ to both $T_{i,j}$ and $T_{j,i}$, and $w_{ij}$ to $X$.
Collect terms, using $w_{ji}=w_{ij}$.
\end{proof}

\paragraph{The types and the sign conditions.}
Since $M_1=0$, nine types remain: $M_2$, $M_3$, the six $T_{i,j}$ and $X$.
All but $X$ are bounded above by Lemma~\ref{lem:corners} and Corollary~\ref{cor:handshake}:
\[
3M_i\le\sum_v\lfloor C_{ii}d_i(v)/2\rfloor\quad\text{for }i=2,3\qquad\text{and}\qquad T_{i,j}\le\sum_v\lfloor C_{ij}d_i(v)/2\rfloor\quad\text{for }i\ne j.
\]
The type $X$ is bounded only below, by $0$.

We impose the \emph{sign conditions}:
\begin{itemize}
\item $w_{12}+w_{13}+w_{23}=0$;
\item $w_2,w_3\ge 0$; and
\item $w_i+2w_{ij}\ge 0$ for all $i\ne j$.
\end{itemize}
A negative coefficient of $X$ would also do, but we do not need one.
In Lemma~\ref{lem:weighted} the sign conditions make the coefficient of $X$ vanish and those of $M_2$, $M_3$ and the six $T_{i,j}$ nonnegative.
The coefficient $3w_1$ of $M_1=0$ may have either sign.

\paragraph{The function $\Phi$.}
Multiplying the bounds on $3M_2$, $3M_3$ and $T_{i,j}$ by $w_2$, $w_3$ and $w_i+2w_{ij}$ then gives $\sum_v\Phi(\mathbf d(v))\le 0$, step~(ii) of the method, for
\[
\Phi(\mathbf d)=W(\mathbf d)-\sum_{i=2,3}w_i\Big\lfloor\frac{C_{ii}d_i}{2}\Big\rfloor-\sum_{i\ne j}(w_i+2w_{ij})\Big\lfloor\frac{C_{ij}d_i}{2}\Big\rfloor .
\]
Dropping the floors only lowers $\Phi$, so the inequality also holds for the floor-free form, which the hand proofs use.
Types with coefficient zero drop out of $\Phi$ and need no cap.

Step~(iii) needs weights for which $\Phi$ is positive at every admissible profile.
At each profile $\Phi$ is linear in the weights, so these requirements and the sign conditions form a finite system of linear inequalities in the six weights.
We give explicit solutions, and the proofs check them directly.

\begin{theorem}\label{thm:345}
$\R(3,4,5)\le 155$.
\end{theorem}

\begin{proof}
\emph{The colour-degree caps.}
Suppose a $(3,4,5;155)$-colouring exists, and write $(r,b,g)=\mathbf d(v)$ for a vertex $v$, so that $r+b+g=154$.
By Lemma~\ref{lem:caps},
\[
r\le D_1=\R(2,4,5)-1=\R(4,5)-1=24\qquad\text{and}\qquad b\le D_2=\R(3,3,5)-1\le 56.
\]

\emph{Choosing which types vanish.}
We choose the weights so that the coefficients of $X$, $T_{1,3}$, $T_{2,1}$ and $T_{3,1}$ in Lemma~\ref{lem:weighted} vanish.
Writing $w_{12}=-\alpha$ and $w_{13}=-\gamma$, this forces $w_1=w_3=2\gamma$, $w_2=2\alpha$ and $w_{23}=\alpha+\gamma$, so that
\[
W=2\gamma\tbinom r2+2\alpha\tbinom b2+2\gamma\tbinom g2-\alpha rb-\gamma rg+(\alpha+\gamma)bg .
\]
As $M_1=0$, Lemma~\ref{lem:weighted} then gives
\[
\sum_vW=2(\gamma-\alpha)T_{1,2}+(4\alpha+2\gamma)T_{2,3}+(2\alpha+4\gamma)T_{3,2}+2\alpha\cdot 3M_2+2\gamma\cdot 3M_3 .
\]
The sign conditions hold exactly when $\gamma\ge\alpha\ge 0$.

\emph{Capping the surviving types.}
By Lemma~\ref{lem:caps}, the surviving types have caps
\begin{gather*}
C_{12}=\R(2,3,5)-1=13\qquad C_{22}=\R(3,2,5)-1=13\\
\text{and}\qquad C_{23}=C_{32}=C_{33}=\R(3,3,4)-1=29.
\end{gather*}
So Corollary~\ref{cor:handshake} and Lemma~\ref{lem:corners} give, without floors,
\[
\begin{gathered}
T_{1,2}\le\tfrac{13}2\textstyle\sum_v r,\quad T_{2,3}\le\tfrac{29}2\sum_v b,\quad T_{3,2}\le\tfrac{29}2\sum_v g,\mkbreak{\quad} 3M_2\le\tfrac{13}2\sum_v b,\quad 3M_3\le\tfrac{29}2\sum_v g.
\end{gathered}
\]
Multiplying these by the coefficients above gives $\sum_v\Phi(\mathbf d(v))\le 0$ for the floor-free form of $\Phi$, here
\[
\Phi=W-13(\gamma-\alpha)r-(71\alpha+29\gamma)b-(29\alpha+87\gamma)g .
\]

\emph{Reduction to two variables.}
Put $x=24-r\ge0$ and $y=56-b\ge 0$, so that $(r,b,g)=(24-x,\,56-y,\,74+x+y)$.
Then
\[
\Phi=(70\alpha-52\gamma)+(70\alpha+132\gamma)x+3\gamma x^2+(47\gamma-63\alpha)y+2(\gamma-\alpha)xy .
\]

Indeed, its three parts are found as follows.
\begin{itemize}
\item \emph{Constant term.}
At $(24,56,74)$ we have $W=5880\alpha+8322\gamma$, and the subtracted terms total $5810\alpha+8374\gamma$.
\item \emph{Linear part.}
The partial derivatives of $\Phi$ there in $r$, $b$, $g$ are $-43\alpha-40\gamma$, $90\alpha+45\gamma$ and $27\alpha+92\gamma$, as $\binom d2$ has derivative $d-\frac12$.
So the linear part is $(70\alpha+132\gamma)x+(47\gamma-63\alpha)y$.
\item \emph{Quadratic part.}
The quadratic part $\gamma r^2+\alpha b^2+\gamma g^2-\alpha rb-\gamma rg+(\alpha+\gamma)bg$ of $W$, at the displacement $(-x,-y,x+y)$, is $3\gamma x^2+2(\gamma-\alpha)xy$.
\end{itemize}

\emph{Fixing the weights.}
We take $\alpha=29$ and $\gamma=39$, that is,
\[
(w_1,w_2,w_3;w_{12},w_{13},w_{23})=(78,58,78;-29,-39,68).
\]
The choice is explained after the proof.
Then
\[
\sum_vW=20\,T_{1,2}+194\,T_{2,3}+214\,T_{3,2}+58\cdot 3M_2+78\cdot 3M_3,
\]
so the sign conditions hold.
Here $\Phi=W-130r-3190b-4234g$, and
\[
\Phi=2+7178x+117x^2+6y+20xy\ \ge\ 2 .
\]
So $\sum_v\Phi(\mathbf d(v))\ge 2\cdot 155>0$, a contradiction.
\end{proof}

\paragraph{The choice of $\alpha$ and $\gamma$.}
Given $\gamma\ge\alpha\ge 0$, the expression for $\Phi$ in the proof is positive on the whole quadrant $x,y\ge 0$ exactly when $70\alpha>52\gamma$ and $47\gamma\ge 63\alpha$, that is,
\[
63/47\le\gamma/\alpha<35/26 .
\]
This interval has width $7/1222$, which explains the thin margin in the proof.
The fraction in it with the least denominator is $39/29$.
With the green cap $g\le\R(3,4,4)-1\le 76$, that is $x+y\le 2$, the lower end would come instead from $\Phi(0,2)=14(3\gamma-4\alpha)>0$.
This would widen the interval only to $4/3<\gamma/\alpha<35/26$, of width $1/78$.

\paragraph{The endpoint $\gamma/\alpha=4/3$.}
The proof of Theorem~\ref{thm:345} needs neither the green cap nor a divisibility argument like parity.
The weights at the endpoint $\gamma/\alpha=4/3$, which the interval above leaves out, also prove it, but only with both.
For $\alpha=3$ and $\gamma=4$,
\[
\Phi=2+738x+12x^2-y+2xy,
\]
which is $-2$ at $(x,y)=(0,4)$, with or without the floors, as $r$, $b$ and $g$ are then even.
Under the green cap $x+y\le 2$, however:
\begin{itemize}
\item if $x=0$, then $\Phi=2-y\ge 0$;
\item if $x\ge 1$, then $\Phi\ge 752$, as $y(2x-1)\ge 0$.
\end{itemize}
So $\Phi\ge 0$, with equality only at $(x,y)=(0,2)$, and $\sum_v\Phi(\mathbf d(v))\le 0$ forces every vertex to have $(r,b,g)=(24,54,76)$.

A divisibility argument modulo~$3$, an analogue of the parity argument, excludes this.
For a vertex $v$ and $u\in N_3(v)$, Lemma~\ref{lem:caps} caps the numbers of red, blue and green neighbours of $u$ in $N_3(v)$ by
\begin{gather*}
C_{31}=\R(2,4,4)-1=\R(4,4)-1=17\\
C_{32}=\R(3,3,4)-1=29\qquad\text{and}\qquad C_{33}=\R(3,4,3)-1=29.
\end{gather*}
These sum to $75=d_3(v)-1$, so all are attained and $e_3(N_3(v))=76\cdot 29/2=1102$.
So $3M_3=155\cdot1102=170810$ by Lemma~\ref{lem:corners}, which is not divisible by~$3$, a contradiction.

\paragraph{The weights of $L$.}
For Theorems~\ref{thm:346} and~\ref{thm:355}, and for $\R(3,5,6)$ in Proposition~\ref{prop:supp}, we require the coefficients of $X$, $T_{1,2}$, $T_{1,3}$, $T_{2,3}$ and $T_{3,2}$ to vanish.
This forces
\[
w_{12}=w_{13}=-w_1/2\qquad w_{23}=w_1\qquad\text{and}\qquad w_2=w_3=-2w_1,
\]
which fixes the weights up to a common factor.
We take
\[
(w_1,w_2,w_3;w_{12},w_{13},w_{23})=(-2,4,4;1,1,-2).
\]
With $(r,b,g)=\mathbf d(v)$, write $L$ for the resulting function $W$.
Then Lemma~\ref{lem:weighted} and $M_1=0$ give
\begin{equation}\label{eq:L}
\begin{gathered}
L=-2\tbinom r2+4\tbinom b2+4\tbinom g2+rb+rg-2bg,\\
\sum_vL(\mathbf d(v))=6T_{2,1}+6T_{3,1}+12M_2+12M_3 .
\end{gathered}
\end{equation}

The remaining coefficients $6=w_2+2w_{12}=w_3+2w_{13}$ and $12=3w_2=3w_3$ are positive, so the sign conditions hold.
The four surviving types have caps $C_{21}$, $C_{22}$, $C_{31}$ and $C_{33}$.
So the terms subtracted from $L$ involve only $b$ and $g$, and the blue-green cap $C_{23}$ used in Theorem~\ref{thm:345} is not needed.

Each floor-free $\Phi$ below is $L$ minus a linear function.
So its quadratic part at a displacement $\delta=(\delta_r,\delta_b,\delta_g)$ is that of $L$, namely
\[
-\delta_r^2+2\delta_b^2+2\delta_g^2+\delta_r\delta_b+\delta_r\delta_g-2\delta_b\delta_g .
\]
Between profiles with the same sum, $\delta_r=-\delta_b-\delta_g$, and this equals $-6\delta_b\delta_g$.
Hence
\[
\Phi(\mathbf d+\delta)=\Phi(\mathbf d)+\nabla\Phi(\mathbf d)\cdot\delta-6\,\delta_b\delta_g\qquad\text{if }\delta_r+\delta_b+\delta_g=0,
\]
where $\nabla\Phi$ is computed with $d-\frac12$ as the derivative of $\binom d2$.
So each proof below needs only a value and a gradient of $\Phi$.

\begin{theorem}\label{thm:346}
$\R(3,4,6)\le 276$.
\end{theorem}

\begin{proof}
\emph{The colour-degree caps.}
Suppose a $(3,4,6;276)$-colouring exists, and write $(r,b,g)=\mathbf d(v)$, so that $r+b+g=275$.
By Lemma~\ref{lem:caps}:
\begin{itemize}
\item $r\le D_1=\R(2,4,6)-1=\R(4,6)-1\le 40$, by the published bound $\R(4,6)\le 41$ rather than the starred $\R(4,6)\le 40$; and
\item $b\le D_2=\R(3,3,6)-1\le 87$ by Theorem~\ref{thm:336}.
\end{itemize}
No cap on $g$ is needed, so Theorem~\ref{thm:345} is not used.

\emph{Capping the surviving types.}
The four types in \eqref{eq:L} have caps
\begin{gather*}
C_{21}=\R(2,3,6)-1=\R(3,6)-1=17\qquad C_{22}=\R(3,2,6)-1=17\\
C_{31}=\R(2,4,5)-1=\R(4,5)-1=24\qquad\text{and}\qquad C_{33}=\R(3,4,4)-1\le 76.
\end{gather*}
So Corollary~\ref{cor:handshake} and Lemma~\ref{lem:corners}, with the floors dropped, give
\begin{gather*}
T_{2,1}\le\frac{17}2\sum_v b\qquad 3M_2\le\frac{17}2\sum_v b\\
T_{3,1}\le 12\sum_v g\qquad\text{and}\qquad 3M_3\le 38\sum_v g.
\end{gather*}
Multiplying these by their coefficients $6$, $4$, $6$ and $4$ in \eqref{eq:L}, where $12M_i=4\cdot 3M_i$, gives $\sum_v\Phi(\mathbf d(v))\le 0$ for the floor-free form of $\Phi$, here
\[
\Phi=L-85b-224g,
\]
with $85=6\cdot\frac{17}2+4\cdot\frac{17}2$ and $224=6\cdot 12+4\cdot 38$.

\emph{Reduction to two variables.}
Put $x=40-r\ge 0$ and $y=87-b\ge 0$, so that $g=148+x+y$.
At $(r,b,g)=(40,87,148)$, $L=40564$ and $85b+224g=40547$, so $\Phi=17$, and $\nabla\Phi=(156,5,232)$.
The expansion after \eqref{eq:L} with $\delta=(-x,-y,x+y)$ gives
\[
\Phi=17+76x+227y+6xy+6y^2\ \ge\ 17 .
\]
So $\sum_v\Phi(\mathbf d(v))\ge 17\cdot 276>0$, a contradiction.
\end{proof}

\paragraph{Why it stops here.}
With the starred $\R(4,6)\le 40$, the cap $r\le 39$ gives $x\ge 1$ in the proof, and the minimum of $\Phi$ rises to $17+76=93$.
Even so, these weights do not exclude $n=275$.
The admissible profile $(39,87,148)$ is $(40,87,148)$ with $r$ lowered by $1$, so $\Phi=17-156-1=-140$ ($-1$ being the coefficient of $r^2$ in $L$).
With the floors it is $-135$, as $17b$ is odd.

\begin{theorem}\label{thm:355}
$\R(3,5,5)\le 346$.
\end{theorem}

\begin{proof}
\emph{The caps.}
Suppose a $(3,5,5;346)$-colouring exists, and write $(r,b,g)=\mathbf d(v)$, so that $r+b+g=345$.
By Lemma~\ref{lem:caps}, $r\le D_1=\R(2,5,5)-1=\R(5,5)-1\le 45$ by $\R(5,5)\le 46$.
As $t_2=t_3$, blue and green play symmetric roles:
\begin{itemize}
\item $b,g\le D_2=D_3=\R(3,4,5)-1\le 154$ by Theorem~\ref{thm:345}; and
\item the four types in \eqref{eq:L} have caps $C_{21}=C_{31}=\R(2,4,5)-1=\R(4,5)-1=24$ and $C_{22}=C_{33}=\R(3,3,5)-1\le 56$.
\end{itemize}

\emph{Capping the surviving types.}
As in Theorem~\ref{thm:346}, without floors,
\[
T_{i,1}\le 12\sum_v d_i(v)\qquad\text{and}\qquad 3M_i\le 28\sum_v d_i(v)\qquad\text{for }i=2,3.
\]
Multiplying these by $6$ and $4$, respectively, gives $\sum_v\Phi(\mathbf d(v))\le 0$ for the floor-free form of $\Phi$, here
\[
\Phi=L-184(b+g),
\]
with $184=6\cdot 12+4\cdot 28$.

\emph{Reduction to two variables.}
Put $x=154-b\ge 0$ and $y=154-g\ge 0$, so that $r=37+x+y$, and $r\le 45$ gives $x+y\le 8$.
At $(r,b,g)=(37,154,154)$, $L=56880=184(b+g)+208$, so $\Phi=208$, and $\nabla\Phi=(235,159,159)$.
The expansion after \eqref{eq:L} with $\delta=(x+y,-x,-y)$ gives
\[
\Phi=208+76x+76y-6xy .
\]
For fixed $y\le 8$, $\Phi$ is increasing in $x$, as $76-6y\ge 28$, so $\Phi\ge 208+76y\ge 208$.
Hence $\sum_v\Phi(\mathbf d(v))\ge 208\cdot 346>0$, a contradiction.
\end{proof}

\paragraph{Why it stops here.}
These weights do not exclude $n=345$: the admissible profile $(36,154,154)$ gives $\Phi=-28$, with or without the floors, as the caps $24$ and $56$ are even.

\paragraph{The gain over the recurrence.}
For comparison, \eqref{eq:recurrence} with our hand bounds $88$ and $155$ in place of $91$ and $157$ gives
\begin{gather*}
41+88+155-1=283\quad\text{for }\R(3,4,6)\\
46+155+155-1=355\quad\text{for }\R(3,5,5).
\end{gather*}
Both are odd, so parity does not apply.
So from the same colour-degree caps counting gains $7$ and $9$, and Theorem~\ref{thm:346} needs only the red and blue ones.
With the starred $\R(4,6)\le 40$, the recurrence gives $282$, lowered by parity to $281$ as $282$ and $88$ are even, and the gain is $5$.

\begin{proposition}\label{prop:supp}
$\R(3,4,7)\le 456$, $\R(3,4,8)\le 700$ and $\R(3,5,6)\le 683$.
Without the starred inputs of Table~\ref{tab:inputs}, $\R(3,4,7)\le 458$, $\R(3,4,8)\le 705$ and $\R(3,5,6)\le 685$.
\end{proposition}

\begin{proof}
\emph{The weights.}
For $\R(3,5,6)$ we use the weights of $L$.
For $\bt=(3,4,k)$ with $k=7,8$ we use
\[
(w_1,w_2,w_3;w_{12},w_{13},w_{23})=(2,2,2;2,-1,-1).
\]
Here $w_{12}+w_{13}+w_{23}=0$ and $w_i+2w_{i3}=w_3+2w_{3i}=0$ for $i=1,2$, while $w_1+2w_{12}=w_2+2w_{12}=3w_2=3w_3=6$.
So Lemma~\ref{lem:weighted} and $M_1=0$ give
\begin{gather*}
W=2\tbinom r2+2\tbinom b2+2\tbinom g2+2rb-rg-bg,\\
\sum_vW(\mathbf d(v))=6T_{1,2}+6T_{2,1}+6M_2+6M_3 .
\end{gather*}

\emph{The caps and the function $\Phi$.}
Of these four types only $M_3$ is counted in green neighbourhoods, with cap $C_{33}=\R(3,4,k-2)-1$.
The other three have the common cap
\[
C:=C_{12}=C_{21}=C_{22}=\R(3,k)-1.
\]
With the weights of $L$, the type $T_{3,1}$ would take the place of $T_{1,2}$, bringing in the larger cap $C_{31}=\R(4,k-1)-1$ multiplied by the green degree $g$.
The bounds on $T_{1,2}$, $T_{2,1}$, $3M_2$ and $3M_3$ enter $\Phi$ with multipliers $6$, $6$, $2$ and $2$, so here
\[
\Phi=W-6\Big\lfloor\frac{Cr}{2}\Big\rfloor-8\Big\lfloor\frac{Cb}{2}\Big\rfloor-2\Big\lfloor\frac{C_{33}g}{2}\Big\rfloor .
\]

\emph{The margins.}
The table gives, for each target:
\begin{itemize}
\item the caps;
\item the number of admissible profiles;
\item the \emph{margin}, in this proof the least value of $\Phi$, with the floors, over these profiles; and
\item a profile attaining it.
\end{itemize}
\begin{center}\small
\begin{tabular}{lcccccc}
\toprule
target & $n$ & $(D_1,D_2,D_3)$ & mixed caps & profiles & margin & attained at\\
\midrule
$(3,4,7)$ & 456 & $(57^*,128,275)$ & $(22,154)$ & 21 & 114 & $(57,128,270)$\\
$(3,4,8)$ & 700 & $(78^*,178,455^*)$ & $(27,275)$ & 91 & 312 & $(78,178,443)$\\
$(3,5,6)$ & 683 & $(84^*,275,345)$ & $(39^*,87,45,154)$ & 276 & 116 & $(62,275,345)$\\
\bottomrule
\end{tabular}
\end{center}

The mixed caps are $(C,C_{33})$ for $(3,4,k)$, and $(C_{21},C_{22},C_{31},C_{33})$ for $(3,5,6)$, that is, $\R(4,6)-1$, $\R(3,3,6)-1$, $\R(5,5)-1$ and $\R(3,4,5)-1$.
A $^*$ marks a cap that rests on a starred input: $\R(4,7)\le 58$, $\R(4,8)\le 79$ and $\R(5,6)\le 85$ in $D_1$, and $\R(4,6)\le 40$ in $C_{21}$.
The cap $D_3=455$ for $(3,4,8)$ comes from the first bound of the proposition, and so it also rests on $\R(4,7)\le 58$.
The other caps come from unstarred entries of Table~\ref{tab:inputs} and the theorems above.

\emph{The margins by hand.}
Every margin is positive, as step~(iii) requires.
The margins can be found by hand, as follows.
Dropping the floors only lowers $\Phi$, and the floor-free minima are already positive.
For the weights $(2,2,2;2,-1,-1)$ write $s=r+b$.
Then
\begin{gather*}
W=s(s-1)+g(g-1)-gs\quad\text{with}\quad g=n-1-s\\
3Cr+4Cb=3Cs+CD_2-C(D_2-b).
\end{gather*}

\begin{itemize}[leftmargin=1.2em]
\item $(3,4,7)$: as $C=22$ and $C_{33}=154$ are even, the floors have no effect, and
\[
\Phi=W-66r-88b-154g=3s^2-1277s+133684+22(128-b),
\]
where $180\le s\le 185$ and $b\le 128$ by the caps.
As $1277/6>185$, the quadratic in $s$ decreases on this range, so the minimum is $114$, at $(r,b,g)=(57,128,270)$.

\item $(3,4,8)$: with $C=27$ and $C_{33}=275$,
\[
\Phi\ge W-81r-108b-275g=3s^2-1903s+290871+27(178-b),
\]
where $244\le s\le 256$ and $b\le 178$.
As $1903/6>256$, the minimum is $311$, at $(78,178,443)$.
There $27\cdot 78$ and $27\cdot 178$ are even and $275\cdot 443$ is odd, so $\Phi$ itself, with the floors, is $312$.
Every other admissible profile has a larger floor-free value, so the margin is $312$.

\item $(3,5,6)$: as in Theorem~\ref{thm:346}, the caps in the table give
\[
\Phi=L-6\Big\lfloor\frac{39b}{2}\Big\rfloor-4\Big\lfloor\frac{87b}{2}\Big\rfloor-6\Big\lfloor\frac{45g}{2}\Big\rfloor-4\Big\lfloor\frac{154g}{2}\Big\rfloor\ \ge\ L-291b-443g,
\]
with $291=6\cdot\frac{39}2+4\cdot\frac{87}2$ and $443=6\cdot\frac{45}2+4\cdot\frac{154}2$.
Put $b=275-x$ and $g=345-y$.
Then $r=62+x+y$, and $r\le 84$ gives $x+y\le 22$.
At $(62,275,345)$, $L-291b-443g$ is $108$ with gradient $(497,179,447)$, so the expansion after \eqref{eq:L} with $\delta=(x+y,-x,-y)$ gives
\[
L-291b-443g=108+318x+50y-6xy.
\]

For fixed $y\le 22$ this increases in $x$, as $318-6y\ge 186$, so it is at least $108+50y\ge 108$, with equality only at $(62,275,345)$.
There the floors add $3$, $2$ and $3$, as $39\cdot 275$, $87\cdot 275$ and $45\cdot 345$ are odd ($154\cdot 345$ is even), so $\Phi$ itself is $116$.
Every other admissible profile has floor-free value at least $158$, so the margin is $116$.
\end{itemize}

\emph{Without the starred inputs.}
For the bounds without starred inputs, use instead:
\begin{itemize}
\item the published $\R(4,7)\le 61$, $\R(4,8)\le 84$ and $\R(5,6)\le 87$ in $D_1$;
\item $\R(4,6)\le 41$ in $C_{21}$; and
\item for $D_3$ of $(3,4,8)$, $\R(3,4,7)\le 458$.
\end{itemize}

The reductions above then give the following.
\begin{itemize}[leftmargin=1.2em]
\item $(3,4,7)$, $n=458$: the caps are now $(60,128,275)$ and $(22,154)$, the floors have no effect, and
\[
\Phi=3s^2-1283s+135198+22(128-b)
\]
on $182\le s\le 188$.
As $1283/6>188$, the minimum is $26$, at $(60,128,269)$, and the margin is $26$.

\item $(3,4,8)$, $n=705$: the caps are now $(83,178,457)$ and $(27,275)$, and
\[
\Phi\ge 3s^2-1918s+296506+27(178-b)
\]
on $247\le s\le 261$.
As $1918/6>261$, the floor-free minimum is $271$, at $(83,178,443)$.
There the floors add $3$, as $27\cdot 83$ is odd, and $1$, as $275\cdot 443$ is odd, so $\Phi=275$.
Every other admissible profile has $s\le 260$ and floor-free value at least $626$, so the margin is $275$.

\item $(3,5,6)$, $n=685$: the caps are now $(86,275,345)$ and $(40,87,45,154)$, and
\[
\Phi\ge L-294b-443g=273+315x+44y-6xy,
\]
with $b=275-x$, $g=345-y$, $r=64+x+y$ and $x+y\le 22$.
As $315-6y\ge 183$, the minimum is $273$, only at $(64,275,345)$.
There the floors add $2$ and $3$, as $87\cdot 275$ and $45\cdot 345$ are odd, so $\Phi=278$.
Every other admissible profile has floor-free value at least $317$, so the margin is $278$.
\end{itemize}

\emph{The computer check.}
A program (Section~\ref{sec:verify}) also evaluates $\Phi$, with the floors, at every admissible profile for both sets of inputs.
It also tests its implementation of Lemma~\ref{lem:weighted}, for both weight vectors, on the $27$ colourings of a triangle.
This suffices, as each corner lies in exactly one triangle, so both sides of the identity are sums over triangles.
\end{proof}

\section{Four-vertex certificates for \texorpdfstring{$\R(4,4,4)$}{R(4,4,4)}}\label{sec:444}

Let $c$ be a $(4,4,4;n)$-colouring.
By Lemma~\ref{lem:caps}, with $\R(3,4,4)\le 77$, $\R(2,4,4)=\R(4,4)=18$ and $\R(3,3,4)=30$, every colour degree is at most $76$, and for $u\in N_i(v)$ we have
\[
|N_i(v)\cap N_i(u)|\le 17\qquad\text{and}\qquad|N_i(v)\cap N_j(u)|\le 29\quad\text{for } j\ne i.
\]

\paragraph{Why triangles do not suffice.}
No colour is now required to be triangle-free, and the weighted triangle identities of Section~\ref{sec:types} give nothing below the recurrence bound $230$.
Their analogue of $\Phi$ is now
\[
W-\sum_iw_i\lfloor 17d_i/2\rfloor-\sum_{i\ne j}(w_i+2w_{ij})\lfloor 29d_i/2\rfloor,
\]
with the sign conditions of Section~\ref{sec:types} and also $w_1\ge 0$, as $M_1$ need not vanish.

Under these conditions, which give rainbow triangles coefficient zero, no weights exclude any order.
Indeed, averaging the weights over the permutations of the colours preserves the conditions and, as the admissible profiles are closed under permuting entries, positivity of $\Phi$ at all of them.
For symmetric weights, however, rainbow triangles have coefficient zero only if $w_{12}=w_{13}=w_{23}=0$, and then
\[
\Phi=w_1\sum_i\big(\binom{d_i}2-\lfloor 17d_i/2\rfloor-2\lfloor 29d_i/2\rfloor\big)
\]
with $w_1\ge 0$.
Each summand is $d_i(d_i-76)/2$, plus $\frac32$ if $d_i$ is odd, and so nonpositive for $0\le d_i\le 76$.

A negative coefficient of $X$, which Section~\ref{sec:types} also allows, does not exclude $n=229$ either.
There every profile is $(76,76,76)$, and $\Phi=1368(w_{12}+w_{13}+w_{23})$ for every choice of weights.
This is $1368$ times the coefficient of $X$, so it is at most $0$ whenever that coefficient is nonpositive.
So we count patterns on four vertices instead.
\begin{itemize}
\item The probabilities that a random star is monochromatic or rainbow depend only on the degree profiles (Lemma~\ref{lem:stars}).
\item The caps bound the probabilities of certain four-vertex events by triangle densities (Lemma~\ref{lem:fivecaps}).
\item Exact identities relate these quantities (Lemma~\ref{lem:444id}).
\end{itemize}
Eliminating the triangle densities gives $\sum_v\Phi(\mathbf d(v))\le 0$ as in step~(ii) (Theorem~\ref{thm:444a}).
Theorem~\ref{thm:444b} uses eight caps, three of them from Theorem~\ref{thm:444a}, and a sum-of-squares term.
Notation introduced in this section is local to it.

\paragraph{The random quadruple.}
Choose an ordered quadruple $(x_0,x_1,x_2,x_3)$ of distinct vertices uniformly at random and, independently, a uniformly random bijection $\sigma$ from the colours $1,2,3$ to labels $R,B,G$.
Write
\[
c_{ij}=\sigma(c(x_ix_j)),
\]
and $N_R(v)$, $N_B(v)$, $N_G(v)$ for the neighbourhoods of $v$ in the colours $R$, $B$, $G$ of $\sigma\circ c$.

The law of $(c_{ij})$ is invariant under permutations of the indices $0,\dots,3$, as the quadruple is uniform, and of the labels, as $\sigma$ is uniform and independent of the quadruple.
So labelled colourings of $K_4$ in one orbit of $S_4\times S_3$ are equally likely.
This holds for every $c$.
No regularity is assumed.

We use the densities
\begin{align*}
s&=P(c_{01}{=}c_{02}{=}c_{03}{=}R), & t&=P(c_{01}{=}R,c_{02}{=}B,c_{03}{=}G),\\
m&=P(c_{01}{=}c_{02}{=}c_{12}{=}R), & b&=P(c_{01}{=}c_{02}{=}R,\,c_{12}{=}B).
\end{align*}
By label symmetry:
\begin{itemize}
\item $3s$ and $6t$ are the probabilities that the star $x_0x_1,x_0x_2,x_0x_3$ is monochromatic, respectively rainbow;
\item $3m$ is the probability that the triangle $x_0x_1x_2$ is monochromatic.
\end{itemize}
Permuting $x_0,x_1,x_2$ and the labels takes the pattern $(c_{01},c_{02},c_{12})=(R,R,B)$ to each of the $18$ patterns with exactly two equal colours.
So each has probability $b$, and $18b$ is the probability that the triangle has exactly two edges of the same colour.

\begin{lemma}[Stars]\label{lem:stars}
Let $v$ be a vertex, $d_i=d_i(v)$ its colour degrees in $c$, and $(x)_3=x(x-1)(x-2)$.
Conditional on $x_0=v$, the star $x_0x_1,x_0x_2,x_0x_3$ is monochromatic with probability $\sum_i(d_i)_3/(n-1)_3$ and rainbow with probability $6d_1d_2d_3/(n-1)_3$.
As $x_0$ is uniform and the two events have unconditional probabilities $3s$ and $6t$, the star densities depend only on the degree profiles:
\[
s=\frac{1}{3n(n-1)_3}\sum_v\sum_i\big(d_i(v)\big)_3,\qquad t=\frac{1}{n(n-1)_3}\sum_v d_1(v)d_2(v)d_3(v).
\]
\end{lemma}

\begin{proof}
Given $x_0=v$, the triple $(x_1,x_2,x_3)$ is uniform among the $(n-1)_3$ ordered triples of distinct vertices other than $v$, and neither event depends on $\sigma$.
The star has all three edges of colour $i$ for $(d_i)_3$ of these triples, and is rainbow for $6d_1d_2d_3$ of them.
Here the factor $6$ counts the ways to assign the three colours to $x_1,x_2,x_3$.
\end{proof}

\paragraph{Pattern probabilities.}
For a pattern $\tau=(c_{01},c_{02},c_{12})$ of the triangle $x_0x_1x_2$ and a condition $\kappa$ on $c_{03}$ and $c_{13}$, let $p(\tau;\kappa)$ be the probability that the triangle has pattern $\tau$ and $\kappa$ holds, and $p(\tau)$ that of $\tau$ alone.
Thus:
\begin{itemize}
\item $p(RRR)=m$;
\item $p(\tau)=b$ for the $18$ patterns $\tau$ with exactly two equal colours, among them $RRB$, $RBR$ and $RBB$;
\item $p(RBG)=(1-3m-18b)/6$, as the six rainbow patterns are equally likely.
\end{itemize}

The table lists the pairs $(\tau,\kappa)$ we use.
In each, $c_{01}=R$, and $\kappa$ says that $x_3$ lies in a set $S$:
\begin{itemize}
\item $S=N_R(x_0)$ if $\kappa$ is $c_{03}=R$;
\item $S=N_R(x_0)\cap N_\chi(x_1)$ if $\kappa$ is $c_{03}=R,\ c_{13}=\chi$ for a label $\chi$.
\end{itemize}
The column ``cap'' bounds $|S|$, and $q$ is the cap less the number of vertices among $x_1,x_2$ that lie in $S$ (Lemma~\ref{lem:fivecaps}).
\begin{center}\small
\begin{tabular}{lcccrr}
\toprule
name & $\tau=(c_{01},c_{02},c_{12})$ & $\kappa$ & cap & $q$ & $\mu$\\
\midrule
$A$ & $RRR$ & $c_{03}=R$ & $76$ & $74$ & \\
$B'$ & $RRR$ & $c_{03}=R,\ c_{13}=B$ & $29$ & $29$ & \\
$\ell_1$ & $RRR$ & $c_{03}=c_{13}=R$ & $17$ & $16$ & $12580$\\
$C$, $\ell_2$ & $RRB$ & $c_{03}=c_{13}=R$ & $17$ & $17$ & $70744$\\
$D$, $\ell_3$ & $RRB$ & $c_{03}=R,\ c_{13}=G$ & $29$ & $29$ & $45584$\\
$\ell_4$ & $RBR$ & $c_{03}=R$ & $76$ & $75$ & $40848$\\
$\ell_5$ & $RBR$ & $c_{03}=R,\ c_{13}=B$ & $29$ & $29$ & $40848$\\
$E$, $\ell_6$ & $RBB$ & $c_{03}=c_{13}=R$ & $17$ & $17$ & $96644$\\
$\ell_7$ & $RBG$ & $c_{03}=R,\ c_{13}=B$ & $29$ & $29$ & $40848$\\
$\ell_8$ & $RBG$ & $c_{03}=R,\ c_{13}=G$ & $29$ & $29$ & $142968$\\
\bottomrule
\end{tabular}
\end{center}

We write $A$, $B'$, $C$, $D$ and $E$ for the probabilities $p(\tau;\kappa)$ of the rows so named.
The prime avoids the label $B$.
For a row marked $\ell_i$, write $\tau_i$, $\kappa_i$, $q_i$, $\mu_i$ for its entries and
\[
\ell_i=p(\tau_i;\kappa_i)-\frac{q_i}{222}\,p(\tau_i).
\]
These are used at $n=225$, where $n-3=222$, in Theorem~\ref{thm:444b}, and $\mu_i$ is the multiplier of $\ell_i$ in \eqref{eq:id225}.

\begin{lemma}[Pattern caps]\label{lem:fivecaps}
In a $(4,4,4;n)$-colouring, with $n'=n-3$, every row of the table satisfies $p(\tau;\kappa)\le\frac{q}{n'}\,p(\tau)$.
In particular,
\[
A\le\tfrac{74}{n'}m,\quad B'\le\tfrac{29}{n'}m,\quad C\le\tfrac{17}{n'}b,\quad D\le\tfrac{29}{n'}b,\quad E\le\tfrac{17}{n'}b .
\]
\end{lemma}

\begin{proof}
Condition on $x_0,x_1,x_2$ and $\sigma$, so that $x_3$ is uniform among the other $n'$ vertices.
If $x_0x_1x_2$ does not have pattern $\tau$, the event has conditional probability $0$.
If it does, then $x_1\in N_R(x_0)$, and Lemma~\ref{lem:caps} with $v=x_0$, and $u=x_1$ for the intersections, gives $|S|\le 76$, $17$ or $29$ according as $S$ is $N_R(x_0)$, $N_R(x_0)\cap N_R(x_1)$ or $N_R(x_0)\cap N_\chi(x_1)$ with $\chi\ne R$.
These caps hold for the labels whatever $\sigma$ is, as the three colours play symmetric roles.

Now:
\begin{itemize}
\item $x_0\notin S$;
\item $x_1\in S$ only for $S=N_R(x_0)$, as $x_1\notin N_\chi(x_1)$;
\item $x_2\in S$ exactly when $\kappa$ holds with $x_2$ in place of $x_3$, which given $\tau$ happens only in rows $A$ and $\ell_1$.
\end{itemize}
So the conditional probability is
\[
|S\setminus\{x_1,x_2\}|/n'\le q/n'.
\]
Averaging over $x_0,x_1,x_2$ and $\sigma$ gives the lemma, and the five inequalities follow as $p(RRR)=m$ and $p(RRB)=p(RBB)=b$.
\end{proof}

\begin{lemma}[Three identities]\label{lem:444id}
For every $3$-colouring of $K_n$ with $n\ge 4$,
\[
\begin{gathered}
\text{(I)}\ \ 9m+18b=1+6s-6t,\mkbreak{\qquad} \text{(II)}\ \ 3A+6B'+6D+6E=2s+m,\qquad \text{(III)}\ \ B'=C.
\end{gathered}
\]
\end{lemma}

\begin{proof}
\emph{Identity (III).}
Exchanging $x_2$ and $x_3$ preserves the law of $(c_{ij})$.
It maps the event of row $C$, $c_{01}=c_{02}=c_{03}=c_{13}=R$ and $c_{12}=B$, onto that of row $B'$, $c_{01}=c_{02}=c_{12}=c_{03}=R$ and $c_{13}=B$.

\emph{Identity (I).}
Both sides equal $3P(c_{01}{=}c_{02})$.
By symmetry in $x_0,x_1,x_2$, this is the expected number of corners of the triangle $x_0x_1x_2$ whose two edges have the same colour:
\begin{itemize}
\item $3$ if the triangle is monochromatic (probability $3m$);
\item $1$ if it has exactly two edges of the same colour ($18b$);
\item $0$ otherwise.
\end{itemize}
By symmetry in $x_1,x_2,x_3$, it is also the expected number of pairs of edges of the star at $x_0$ with the same colour:
\begin{itemize}
\item $3$ if the star is monochromatic ($3s$);
\item $0$ if it is rainbow ($6t$);
\item $1$ otherwise ($1-3s-6t$).
\end{itemize}
So
\[
9m+18b=9s+(1-3s-6t)=1+6s-6t.
\]

\emph{Identity (II): reduction.}
Split $s$ by $c_{12}$ and $m$ by $c_{03}$.
Exchanging the labels $B$ and $G$ preserves the law, so the cases $c_{12}=G$ of $s$ and $c_{03}=G$ of $m$ have the same probabilities as $c_{12}=B$ and $c_{03}=B$.
The case $c_{12}=R$ of $s$ and the case $c_{03}=R$ of $m$ are both the event of row $A$.
Splitting further by $c_{13}$:
\begin{itemize}
\item the case $c_{12}=B$ of $s$ gives $C$, $D$ and $Y_s$ for $c_{13}=R$, $G$, $B$;
\item the case $c_{03}=B$ of $m$ gives $B'$, $E$ and $Y_m$ for $c_{13}=R$, $B$, $G$, the values $B'$ and $E$ after exchanging $x_0$ with $x_1$, respectively $x_2$ with $x_3$.
\end{itemize}
Here
\[
Y_s=P(c_{01}{=}c_{02}{=}c_{03}{=}R,\ c_{12}{=}c_{13}{=}B),\qquad Y_m=P(c_{01}{=}c_{02}{=}c_{12}{=}R,\ c_{03}{=}B,\ c_{13}{=}G).
\]
So
\[
s=A+2(C+D+Y_s)\qquad\text{and}\qquad m=A+2(B'+E+Y_m),
\]
and with (III),
\[
2s+m=3A+6B'+4D+2E+4Y_s+2Y_m.
\]
Hence (II) is equivalent to
\[
2Y_s+Y_m=D+2E.
\]

\emph{Identity (II): the twelve terms.}
To prove this, split each of $Y_s$, $Y_m$, $D$ and $E$ further by $c_{23}$.
Each of the twelve resulting terms is the probability of one labelled colouring of $K_4$, and terms whose colourings lie in one $S_4\times S_3$-orbit are equal.
The table gives, for each term, the classes of colours $R$, $B$ and $G$, in this order, as graphs on $\{x_0,\dots,x_3\}$.
Here the paw is a triangle with a pendant edge, $P_3$ the path with two edges, and $\emptyset$ an empty class.
\begin{center}\small
\begin{tabular}{lccc}
\toprule
 & $c_{23}=R$ & $c_{23}=B$ & $c_{23}=G$\\
\midrule
$Y_s$ & paw, $P_3$, $\emptyset$\ \ ($\pi_1$) & $K_{1,3}$, $K_3$, $\emptyset$\ \ ($\pi_2$) & $K_{1,3}$, $P_3$, $K_2$\ \ ($\pi_3$)\\
$Y_m$ & paw, $K_2$, $K_2$\ \ ($\pi_4$) & $K_3$, $P_3$, $K_2$\ \ ($\pi_5$) & $K_3$, $K_2$, $P_3$\ \ ($\pi_5$)\\
$D$ & paw, $K_2$, $K_2$\ \ ($\pi_4$) & $K_{1,3}$, $P_3$, $K_2$\ \ ($\pi_3$) & $K_{1,3}$, $K_2$, $P_3$\ \ ($\pi_3$)\\
$E$ & paw, $P_3$, $\emptyset$\ \ ($\pi_1$) & $K_3$, $K_{1,3}$, $\emptyset$\ \ ($\pi_2$) & $K_3$, $P_3$, $K_2$\ \ ($\pi_5$)\\
\bottomrule
\end{tabular}
\end{center}

\emph{Identity (II): the orbits.}
Colourings with the same multiset of classes lie in one orbit.
Indeed, each multiset contains a paw, $K_{1,3}$ or $K_3$.
Permuting labels and vertices, we may assume that both colourings have this class, in the same colour.
Its stabiliser in $S_4$, together with the exchange of the other two labels, acts transitively on the splittings of its complement into the remaining two classes.
Writing $\pi_1,\dots,\pi_5$ for the five common values, as marked, the rows give
\[
Y_s=\pi_1+\pi_2+\pi_3,\quad Y_m=\pi_4+2\pi_5,\quad D=\pi_4+2\pi_3,\quad E=\pi_1+\pi_2+\pi_5,
\]
and hence
\[
2Y_s+Y_m=2\pi_1+2\pi_2+2\pi_3+\pi_4+2\pi_5=D+2E.
\]
\end{proof}

\begin{theorem}\label{thm:444a}
$\R(4,4,4)\le 226$.
\end{theorem}

\begin{proof}
\emph{Step 1: combining the caps.}
Suppose that a $(4,4,4;226)$-colouring exists, so that $n'=223$.
For $0\le\lambda\le 1$, add the five inequalities of Lemma~\ref{lem:fivecaps} with multipliers $\frac12$, $\lambda$, $1-\lambda$, $1$ and $1$.
The multipliers $\lambda$ and $1-\lambda$ divide the one probability $B'=C$ between its two caps.
By (III) and (II), the left-hand side is $\frac12A+B'+D+E=\frac16(2s+m)$, so
\[
\tfrac16(2s+m)\le\frac{(37+29\lambda)m+(63-17\lambda)b}{n'} .
\]

\emph{Step 2: eliminating $m$.}
Substituting $b=(1+6s-6t)/18-m/2$ from (I) gives $m$ the coefficient $(11+75\lambda)/(2n')$ on the right.
If this is at most $\frac16$, the net term in $m$ on the left is nonnegative, as $m\ge 0$, and can be dropped, leaving an inequality in $s$ and $t$.
As $1+6s-6t=9m+18b\ge 0$, this is strongest when the coefficient $63-17\lambda$ of $b$ is least.
So we take the largest such $\lambda$, namely $\lambda=(n'-33)/225$, here $\frac{38}{45}$, for which the terms in $m$ cancel.

\emph{Step 3: the star inequality.}
Multiplying by $18n'$ then gives $6n's\le(63-17\lambda)(1+6s-6t)$, and substituting $\lambda=(n'-33)/225$ and multiplying by $225$ gives
\[
1350\,n's\le(14736-17n')(1+6s-6t).
\]
For $n'=223$, dividing by $5$ gives $60210\,s\le 2189(1+6s-6t)$, that is,
\begin{equation}\label{eq:star226}
15692\cdot 3s+2189\cdot 6t\le 2189,
\end{equation}
or $15692\,P(\text{monochromatic star})+2189\,P(\text{rainbow star})\le 2189$ for the star at $x_0$.

\emph{Step 4: the profiles.}
Every vertex has a profile with entries at most $76$ summing to $225$.
Up to order, these are $(73,76,76)$, $(74,75,76)$ and $(75,75,75)$.
Given $x_0=v$, Lemma~\ref{lem:stars} gives the conditional probabilities $\sum_i(d_i)_3/(n-1)_3$ and $6d_1d_2d_3/(n-1)_3$ of a monochromatic and a rainbow star, where $d_i=d_i(v)$.
With these in place of $3s$ and $6t$, the left side of \eqref{eq:star226} minus $2189$ becomes
\[
\Phi(\mathbf d(v))=\frac{15692\sum_i(d_i)_3+6\cdot 2189\,d_1d_2d_3-2189\,(n-1)_3}{(n-1)_3},
\]
where
\[
(n-1)_3=225\cdot 224\cdot 223=11{,}239{,}200\qquad\text{and}\qquad 2189\,(n-1)_3=24{,}602{,}608{,}800.
\]

\emph{Step 5: the contradiction.}
As $x_0$ is uniform, \eqref{eq:star226} says that
\[
\sum_v\Phi(\mathbf d(v))\le 0.
\]
But $\Phi$ is symmetric in the $d_i$ and positive at every profile:
\begin{center}\small
\begin{tabular}{lrrrr}
\toprule
profile & $\sum_i(d_i)_3$ & $d_1d_2d_3$ & numerator & $\Phi(\mathbf d)$\\
\midrule
$(73,76,76)$ & $1{,}216{,}776$ & $421{,}648$ & $28{,}965{,}024$ & $451/175$\\
$(74,75,76)$ & $1{,}215{,}894$ & $421{,}800$ & $17{,}121{,}048$ & $457/300$\\
$(75,75,75)$ & $1{,}215{,}450$ & $421{,}875$ & $11{,}138{,}850$ & $111/112$\\
\bottomrule
\end{tabular}
\end{center}
This is a contradiction.
\end{proof}

\paragraph{Where the multipliers come from.}
The multipliers were found by linear programming at $n=228$, where every profile is $(75,76,76)$, which selected the five caps used here.
With these caps, (II), (III) and the cancellation of $m$ determine the multipliers up to a common factor.
At $n=226$ exact linear algebra gives $\frac12$, $\frac{38}{45}$, $\frac{7}{45}$, $1$ and $1$.
The proof does not depend on this.

\paragraph{Why it stops here.}
At $n=225$, where $n'=222$ and $\lambda=\frac{21}{25}$, the inequality $1350\,n's\le(14736-17n')(1+6s-6t)$ of the proof, divided by $54$, becomes $5550\,s\le 203(1+6s-6t)$, that is,
\[
1444\cdot 3s+203\cdot 6t\le 203
\]
in place of \eqref{eq:star226}.

The profiles now have entries at most $76$ summing to $224$.
Up to order they are $(72,76,76)$, $(73,75,76)$, $(74,74,76)$ and $(74,75,75)$.
At each, the left side minus $203$, computed as in the proof, is negative.
In this order the values are $-1369/1561$, $-1597/1561$, $-239/223$ and $-28009/24976$, the largest being $-1369/1561\approx-0.88$.
Since the proof takes the best $\lambda$ for which the term in $m$ can be dropped, no choice of $\lambda$ lets this argument exclude $225$.

\paragraph{A sum of squares.}
Theorem~\ref{thm:444b} combines the eight inequalities $\ell_i\le 0$ of the table, which hold in a $(4,4,4;225)$-colouring by Lemma~\ref{lem:fivecaps}, and a sum-of-squares term $Q\ge 0$, through one linear identity checked by computer.
Three of the $\ell_i$ are the caps on $C$, $D$ and $E$.
The caps on $A$ and $B'$ are not used.

For a $3$-colouring of $K_{225}$, define $Q$ as follows.
Identify the labels $R$, $B$, $G$ with the residues $0$, $-1$, $1$ modulo $3$, and let $a(\alpha,\beta)\in\{-1,0,1\}$ be congruent to $\alpha+\beta$.
For each vertex $w\ne x_0,x_1$, let
\[
a_w=a\big(\sigma(c(x_0w)),\sigma(c(x_1w))\big)\qquad\text{and}\qquad Z=\sum_wa_w,
\]
so that $Z$ is a sum of $n-2=223$ terms depending on $x_0$, $x_1$ and $\sigma$.
Then
\[
Q:=\mathbb E\big[\mathbf 1_{c_{01}=R}\,(Z/223)^2\big]=\tfrac{222}{223}\,\mathbb E\big[\mathbf 1_{c_{01}=R}\,a_{x_2}a_{x_3}\big]+\tfrac1{223}\,\mathbb E\big[\mathbf 1_{c_{01}=R}\,a_{x_2}^2\big]\ \ge 0 .
\]

Indeed,
\[
Z^2=\sum_wa_w^2+\sum_{w\ne w'}a_wa_{w'},
\]
where the second sum has $223\cdot 222$ terms.
Given $x_0$, $x_1$ and $\sigma$, the indicator is fixed and $(x_2,x_3)$ is a uniform ordered pair of distinct vertices among the other $223$.
So the conditional means of $a_{x_2}^2$ and $a_{x_2}a_{x_3}$ are the averages of the terms of the two sums.
As $a_{x_j}=a(c_{0j},c_{1j})$, both expectations on the right are combinations of probabilities of colour patterns on $x_0,\dots,x_3$.

Any real function $a$ would give $Q\ge 0$.
This one is needed only for the identity \eqref{eq:id225} below.

\begin{theorem}[identity checked by computer]\label{thm:444b}
$\R(4,4,4)\le 225$.
\end{theorem}

\begin{proof}
\emph{Step 1: the identity.}
Besides $s$ and $t$ we use a third star density,
\[
\eta=P(c_{01}{=}c_{02}{=}R,\ c_{03}{=}B).
\]
As for $b$, each of the $18$ patterns $(c_{01},c_{02},c_{03})$ of the star with exactly two equal colours has probability $\eta$, so $3s+18\eta+6t=1$.
Thus $s$ and $t$ determine $\eta$.
We keep all three so that the left side of \eqref{eq:id225} has no constant term.
For every $3$-colouring of $K_{225}$,
\begin{equation}\label{eq:id225}
\tfrac{35020}{3}\,s+\tfrac{39040}{3}\,\eta-24012\,t=\sum_{i=1}^8\mu_i\ell_i-5129\,Q .
\end{equation}

\emph{Step 2: checking the identity.}
With $Q$ expanded as above, both sides are linear combinations of the probabilities of the $3^6=729$ labelled colourings of $K_4$ on $(x_0,\dots,x_3)$.
By the invariance noted at the start of the section, these probabilities are constant on each of the $15$ orbits of $S_4\times S_3$.
So it suffices to check \eqref{eq:id225} for the uniform law on each orbit.
We did this in exact rational arithmetic, with a program (\texttt{verify\_r444.py}, Section~\ref{sec:verify}) that shares no code with the search that found the multipliers.

\emph{Step 3: the sign of the right side.}
In a $(4,4,4;225)$-colouring the right side of \eqref{eq:id225} is at most $0$, as each $\mu_i$ is positive, each $\ell_i\le 0$ and $Q\ge 0$.

\emph{Step 4: the left side at a vertex.}
Given $x_0=v$, write $d_i=d_i(v)$.
By Lemma~\ref{lem:stars} and label symmetry, the events defining $s$ and $t$ have conditional probabilities $\sum_i(d_i)_3/(3(n-1)_3)$ and $d_1d_2d_3/(n-1)_3$, where $(n-1)_3=224\cdot223\cdot222$.

For $\eta$, the colours of $c$ that $\sigma$ maps to $R$ and $B$ form a uniformly random ordered pair $(i,j)$ of distinct colours.
Given it, $x_0x_1$ and $x_0x_2$ have colour $i$ and $x_0x_3$ colour $j$ for $d_i(d_i-1)d_j$ of the $(n-1)_3$ triples $(x_1,x_2,x_3)$.
So the event defining $\eta$ has conditional probability $\sum_{i\ne j}d_i(d_i-1)d_j/(6(n-1)_3)$.
Multiplying the conditional value of the left side of \eqref{eq:id225} by $18(n-1)_3$ and using $\sum_{j\ne i}d_j=n-1-d_i=224-d_i$ gives
\[
J(\mathbf d)=70040\sum_i(d_i)_3+39040\sum_id_i(d_i-1)(224-d_i)-432216\,d_1d_2d_3 .
\]

\emph{Step 5: the contradiction.}
As $x_0$ is uniform, the left side of \eqref{eq:id225} is the average over $v$ of $J(\mathbf d(v))/(18(n-1)_3)$.
So $\sum_vJ(\mathbf d(v))\le 0$, and $J$ plays the role of $\Phi$.
But $J$ is symmetric, and positive at each of the four profiles listed after Theorem~\ref{thm:444a}:
\begin{center}\small
\begin{tabular}{lrrrr}
\toprule
profile & $\sum_i(d_i)_3$ & $\sum_id_i(d_i-1)(224-d_i)$ & $d_1d_2d_3$ & $J(\mathbf d)$\\
\midrule
$(72,76,76)$ & $1{,}201{,}440$ & $2{,}464{,}224$ & $415{,}872$ & $605{,}630{,}208$\\
$(73,75,76)$ & $1{,}200{,}126$ & $2{,}464{,}206$ & $416{,}100$ & $414{,}349{,}680$\\
$(74,74,76)$ & $1{,}199{,}688$ & $2{,}464{,}200$ & $416{,}176$ & $350{,}589{,}504$\\
$(74,75,75)$ & $1{,}199{,}244$ & $2{,}464{,}200$ & $416{,}250$ & $287{,}507{,}760$\\
\bottomrule
\end{tabular}
\end{center}
This is a contradiction.
\end{proof}

\paragraph{How the certificate was found.}
The function $a$, the $\mu_i$ and $5129$ were found numerically.
The proof does not depend on how.
At $n=225$, consider the linear programme in the probabilities of the $14$ orbits of colourings of $K_4$ that are not monochromatic, subject to the following constraints:
\begin{itemize}
\item the star densities are a mixture of those of the admissible profiles;
\item the degree bounds $72\le d_i\le 76$ and the caps $17$ and $29$ of Lemma~\ref{lem:caps} hold, conditioned on each of the $27$ patterns of $x_0x_1x_2$.
\end{itemize}
This programme is feasible.

At a feasible point, the $10\times 10$ matrix of the quadratic form $(v_0,v)\mapsto\mathbb E[\mathbf 1_{c_{01}=R}(v_0+Z_v/223)^2]$, where $Z_v$ is $Z$ with $a$ replaced by a function $v$ on pairs of labels, has negative eigenvalues.
For an eigenvector of the most negative one, the signs of its nine entries indexed by pairs of labels give $a$.
The single constraint $Q\ge 0$ then makes the programme infeasible, and the $\mu_i$ and $5129$ are a dual solution, rationalised and then checked exactly.
Given the terms of \eqref{eq:id225}, the identity determines its coefficients up to a common factor, so $\mu_4=\mu_5=\mu_7$ is forced.

The residues modulo~$3$ in the definition of $Q$ only encode the nine values of $a$.
They are unrelated to the argument of \cite{Boz26}.

\paragraph{Other values of $n$.}
With the following changes, the same method gives exact certificates for $219\le n\le 224$ (Section~\ref{sec:verify}):
\begin{itemize}
\item several squares $\mathbf 1_{c_{01}=R}(v_0+Z_v/(n-2))^2$ in place of $Q$;
\item three to five caps of the kind in Lemma~\ref{lem:fivecaps};
\item in place of the identity, conditions (i) and (ii) of Lemma~\ref{lem:flagcert} on four vertices, for the star term $c_s\mathbf 1_s+c_\eta\mathbf 1_\eta+c_t\mathbf 1_t$, where $\mathbf 1_s$, $\mathbf 1_\eta$ and $\mathbf 1_t$ are the indicators of the events defining $s$, $\eta$ and $t$, these caps and these squares;
condition (i) is checked on the $14$ orbits other than that of a monochromatic $K_4$, which has probability $0$.
\end{itemize}
Theorem~\ref{thm:r444five} supersedes them.

\section{Five-vertex certificates with local caps}\label{sec:five}

Section~\ref{sec:444} averages over a random ordered quadruple of vertices.
The same idea with a random ordered $5$-tuple, the caps of Lemma~\ref{lem:multicaps} and many squares gives most of our strongest bounds.

The certificates are too large to print.
Instead we describe their form, prove that a certificate of this form excludes the order $n$, and give their sizes and margins.
Every certificate was found in floating point, by semidefinite programming, and then rounded to rational numbers.
The solvers were:
\begin{itemize}
\item for Theorem~\ref{thm:r444five}, a dual log-barrier interior-point method written for this work;
\item for Theorem~\ref{thm:r34five} and the mixture-model certificate for $\R(4,4,4)\le 200$ (see after Theorem~\ref{thm:r444five}), Clarabel \cite{GC24}, with the star and row multipliers re-solved by HiGHS \cite{HH18} after the squares are rounded;

\item for $\R(3,3,k)$, a cutting-plane method for the semidefinite constraints, with the linear programmes solved by HiGHS.
\end{itemize}
Every certificate is checked in exact rational arithmetic by a program that shares no code with the search (Section~\ref{sec:verify}).

\subsection{Certificates}\label{sec:cert}

Let $c$ be a $(\bt;n)$-colouring and $s\in\{4,5\}$.
Let $\Gamma$ be the group of permutations $\gamma$ of the colours with $t_{\gamma(i)}=t_i$ for all $i$.
For $\bt=(3,3,k)$ we use instead the group of order $2$ exchanging red and blue, also when $k=3$; everything below holds for any group of colour permutations fixing $\bt$.
Choose an ordered $s$-tuple $(x_0,\dots,x_{s-1})$ of distinct vertices uniformly at random and, independently, a uniform $\sigma\in\Gamma$.
Let $P$ be the colouring $x_ix_j\mapsto\sigma(c(x_ix_j))$ of $K_s$.
Neighbourhoods below are taken in $\sigma\circ c$, which is again $\bt$-avoiding.

As in Section~\ref{sec:444}, the law of $P$ is invariant under $S_s\times\Gamma$.
So
\[
\mathbb E f(P)=\mathbb E\,(\mathrm{Sym} f)(P)
\]
for every function $f$ on colourings of $K_s$, where $\mathrm{Sym} f$ averages $f$ over the orbit; and $P$ has no colour-$i$ $K_{t_i}$ for any $i$.

We use three kinds of terms.
\begin{enumerate}[label=(\alph*)]
\item \emph{Stars.}
For a pattern $\pi$ of colours on the star $x_0x_1,\dots,x_0x_{s-1}$, the probability that the star has pattern $\pi$, given $x_0=v$, is a polynomial of degree $s-1$ in $\mathbf d(v)$, as in Lemma~\ref{lem:stars}.
So for $G=\sum_\pi y_\pi\mathbf 1[\pi]$ we have
\[
\mathbb E\,G=\mathbb E_v J_G(\mathbf d(v)),
\]
where $v$ is a uniform vertex and $J_G$ is an explicit polynomial.

\item \emph{Rows.}
Fix a colouring $\rho$ of the pairs among $x_0,\dots,x_{s-2}$, a nonempty set $Y$ of these vertices and colours $a(y)$, $y\in Y$, and let
\[
S=\bigcap_{y\in Y}N_{a(y)}(y).
\]
Lemma~\ref{lem:multicaps}, with $\mathbf w$ read off from $\rho$, caps $|S|$.
Let $q$ be this cap less the number of vertices among $x_0,\dots,x_{s-2}$ that lie in $S$, which $\rho$ determines.
\begin{itemize}
\item Given $x_0,\dots,x_{s-2}$, the vertex $x_{s-1}$ is uniform among the other $n-s+1$.
So the \emph{cap row}
\[
g=\mathbf 1[\rho]\big(\mathbf 1[x_{s-1}\in S]-q/(n-s+1)\big)
\]
has $\mathbb E\,g\le 0$.
\item Let $Y=\{y\}$ and $a(y)=i$, so that $S=N_i(y)$, and let $q$ be the lower bound
\[
d_i(y)\ge n-1-\sum_{l\ne i}D_l
\]
less the number of vertices among $x_0,\dots,x_{s-2}$ that lie in $S$.
Then the \emph{lower row}
\[
g=\mathbf 1[\rho]\big(q/(n-s+1)-\mathbf 1[x_{s-1}\in S]\big)
\]
has $\mathbb E\,g\le 0$.
\item For $\bt=(3,3,k)$ we also use \emph{profile rows}: functions $g$ of the colours of $x_0x_1$, $x_0x_2$ and $x_1x_2$ with $\mathbb E\,g\le\mathbb E_v\beta(\mathbf d(v))$ for an explicit function $\beta$.
For fixed $\sigma$, given $x_0=v$, the probability that $x_0x_1$ and $x_0x_2$ have colour $i$ and $x_1x_2$ colour $j$ is $2e_j(N_i(v))/((n-1)(n-2))$, with $N_i$ taken in $\sigma\circ c$.
So bounds on $e_j(N_i(v))$ give three kinds of profile row.
Each $\beta$ below is the bound for fixed $\sigma$, with $d_i$ the colour-$i$ degree of $v$ in $\sigma\circ c$.
The row uses its average over $\sigma\in\Gamma$, that is, over the exchange of red and blue, as a function of $\mathbf d(v)$:
\begin{itemize}
\item the \emph{upper floor row}, the indicator $g$ of this event, with $\beta(\mathbf d)=2\lfloor C_{ij}d_i/2\rfloor/((n-1)(n-2))$, by Corollary~\ref{cor:handshake};
\item the \emph{lower floor row} $-g$, with $\beta(\mathbf d)=-2\big\lceil d_i\max\big(0,d_i-1-\sum_{j'\ne j}C_{ij'}\big)/2\big\rceil/((n-1)(n-2))$, as each $u\in N_i(v)$ has at least $d_i-1-\sum_{j'\ne j}C_{ij'}$ colour-$j$ neighbours in $N_i(v)$;
\item the \emph{recursive row}, minus the indicator that $x_0x_1$ and $x_0x_2$ are green and $x_1x_2$ is red or blue, with $\beta(\mathbf d)=-2m_{k-1}(d_3)/((n-1)(n-2))$, by the valid lower bound $e_H(N_3(v))\ge m_{k-1}(d_3(v))$ from the proof of Lemma~\ref{lem:density}.
\end{itemize}
\end{itemize}
For cap and lower rows put $\beta=0$.

\item \emph{Squares.}
Fix a root, a vertex, an edge or a triangle with given colours, placed on the first one, two or three vertices of the tuple.
Give every vertex $w$ of $K_n$ outside the root (for a vertex root, every pair of such vertices, ordered or unordered) a \emph{type}, the colours of its edges to the root (and between the pair).
For a constant $\theta_0$ and a vector $\theta$ indexed by the types, let
\[
Z=\sum_w\theta_{\mathrm{type}(w)},
\]
summed over these vertices (or pairs), and let $N$ be the number of terms.
Then
\[
\mathbf 1[\text{root}]\,(\theta_0+Z/N)^2\ge 0.
\]
As in the definition of $Q$ in Section~\ref{sec:444}, averaging over the remaining vertices of the tuple turns its expectation into that of a function of $P$, provided the positions of the tuple outside the root can hold two summed vertices (or two disjoint pairs).
Indeed, the square expands into $\theta_0^2$, $2\theta_0$ times a mean over one vertex (or pair) outside the root, and $Z^2/N^2$.
The last splits by how the two summed vertices or pairs overlap: equal or distinct vertices, and pairs sharing two vertices, one or none.
Each part is a mean over at most four vertices outside the root, and so it is the expectation of a function of $P$.
In Lemma~\ref{lem:flagcert}, $Q_j$ denotes this function of $P$, which depends on $n$ through $N$.
The pairs are ordered in the certificates for $\R(3,3,k)$ and those of Theorem~\ref{thm:r34five}, and unordered in those of Theorem~\ref{thm:r444five}.
In the certificates of Theorem~\ref{thm:r34five} a square at a vertex root may also combine the means over single vertices and over pairs (that is, $Z/N$ is replaced by $Z_1/N_1+Z_2/N_2$, with $Z_1$ summed over single vertices and $Z_2$ over ordered pairs).
For $\bt=(4,4,4)$ we also use \emph{localised} squares at a red edge $x_0x_1$:
\begin{itemize}
\item $\mathbf 1[\text{root}]\,(q-|S|)(\theta_0+Z/N)^2\ge 0$, where $S$ is $N_i(x_0)\cap N_j(x_1)$, $N_i(x_0)$ or $N_j(x_1)$ with $x_0,x_1$ removed, and $q$ is the cap of Lemma~\ref{lem:multicaps} less the number of removed vertices that lie in the set; and
\item $\mathbf 1[\text{root}]\,(|S|-q)(\theta_0+Z/N)^2\ge 0$ for $S=N_i(x_0)\setminus\{x_1\}$, with $q$ the lower degree bound less $\mathbf 1[x_1\in N_i(x_0)]$.
\end{itemize}
Writing $|S|=\sum_u\mathbf 1[u\in S]$ over the vertices $u\ne x_0,x_1$, and splitting again by whether $u$ equals one of the two summed vertices, turns the expectation into that of a function of $P$ on at most three vertices outside the root.
\end{enumerate}

\begin{lemma}\label{lem:flagcert}
Suppose that for a star term $G$, rows $g_r$ with $\mathbb E\,g_r\le\mathbb E_v\beta_r(\mathbf d(v))$, squares $Q_j$ (so $\mathbb E\,Q_j\ge 0$) and multipliers $\mu_r,\lambda_j\ge 0$,
\begin{enumerate}[label=(\roman*)]
\item the inequality
\[
\mathrm{Sym}\Big(G-\sum_r\mu_rg_r+\sum_j\lambda_jQ_j\Big)\le 0
\]
holds at every colouring of $K_s$ that has no colour-$i$ $K_{t_i}$ for any $i$, and
\item the inequality
\[
J_G(\mathbf d)-\sum_r\mu_r\beta_r(\mathbf d)>0
\]
holds at every admissible profile $\mathbf d$.
\end{enumerate}
Then no $(\bt;n)$-colouring exists.
If instead of (ii),
\[
(n-1)\Big(J_G(\mathbf d)-\sum_r\mu_r\beta_r(\mathbf d)\Big)+d_1+d_2\ge h^*
\]
at every admissible profile, and $\bt=(3,3,k)$, then $\lceil nh^*/2\rceil$ is a valid lower bound $m_k(n)$.
\end{lemma}

\begin{proof}
By (i) and the invariance of the law of $P$,
\[
0\ge\mathbb E\Big(G-\sum_r\mu_rg_r+\sum_j\lambda_jQ_j\Big)\ge\mathbb E_v\Big(J_G(\mathbf d(v))-\sum_r\mu_r\beta_r(\mathbf d(v))\Big),
\]
which (ii) makes positive, as every profile is admissible.
Under the alternative hypothesis the same chain gives
\[
0\ge\big(h^*-\mathbb E_v(d_1(v)+d_2(v))\big)/(n-1),
\]
so the average $H$-degree is at least $h^*$, and $e(H)\ge nh^*/2$.
\end{proof}

The star term absorbs any constant, as $\sum_\pi\mathbf 1[\pi]=1$.
So (i) and (ii) say that the profile polynomial, corrected by the rows, is positive, while the combination of caps and squares makes the whole function nonpositive.

Theorems~\ref{thm:444a} and~\ref{thm:444b} are the case $s=4$, with no square or one, and with identities in place of the inequality (i).
The certificates for $219\le n\le 224$ described at the end of Section~\ref{sec:444} are the case $s=4$ with several squares.
In the floating-point search described at the start of the section, the vectors of the squares are rounded to integers and the linear part is re-solved before the multipliers are rationalised.

Checking (i) needs only the orbits of $S_s\times\Gamma$ on colourings of $K_s$ without a forbidden clique.
For $s=5$ there are $133$ for $(4,4,4)$, $193$ for $(3,3,k)$ with $k\ge 6$, $273$ for $(3,4,4)$, and between $287$ and $566$ for the other targets.
Checking (ii) needs every admissible profile.

\subsection{Results}

\begin{theorem}[certificate checked by computer]\label{thm:r444five}
$\R(4,4,4)\le 198$.
\end{theorem}

\begin{proof}
A certificate as in Lemma~\ref{lem:flagcert} for $\bt=(4,4,4)$, $n=198$ and $s=5$.
\begin{itemize}
\item Its star term has four classes, the multiplicity patterns $4$, $3{+}1$, $2{+}2$ and $2{+}1{+}1$ of the colours on $x_0x_1,\dots,x_0x_4$, so $J_G$ is a quartic in the profile.
\item It has $1581$ cap and lower rows, conditioned on the colouring of $x_0,\dots,x_3$, with $Y$ of size $1$ to $4$.
\item It has $233$ squares:
\begin{itemize}
\item $113$ at a red edge with $9$ types, at a monochromatic, a two-coloured and a rainbow triangle with $27$ types, and at a vertex with $18$ types of unordered pairs, and
\item $120$ localised squares at a red edge, with the constants $q=17$, $29$, $75$ and $76$ (the cap $76$ less $1$ gives $75$ for $N_1(x_0)$, which contains $x_1$), and with the lower degree bound.
\end{itemize}
\item The inputs are $\R(3,4,4)\le 77$, so that every colour degree is at most $76$, $\R(3,3,4)=30$, $\R(3,3,3)=17$, $\R(4,4)=18$, $\R(3,4)=9$ and $\R(3,3)=6$.
\item The multipliers are dyadic rationals and the square vectors are integral.
\end{itemize}
In exact arithmetic, the largest value in (i) over the $133$ orbits is $-3.05\cdot 10^{-9}$, and the least value in (ii) over the $528$ admissible profiles is $3.05\cdot 10^{-9}$.
A second program views every term of the certificate as a function of the colouring of $K_6$ on $(x_0,\dots,x_5)$, which changes no expectation, as no term involves $x_5$.
It evaluates (i) over the $3929$ orbits of $S_6\times S_3$ on colourings of $K_6$ without a monochromatic $K_4$, by exact integer counts over all $3^{15}$ labelled colourings, and finds the largest value in (i) to be $-3.05\cdot 10^{-9}$ as well (to these digits; exactly, it is slightly smaller) and the same least value in (ii).
\end{proof}

\paragraph{Comparison with other versions.}
Theorem~\ref{thm:r444five} supersedes Theorems~\ref{thm:444a} and~\ref{thm:444b}.
Below, the relaxation at order $n$ is \emph{decided} if the floating-point optimum of its semidefinite programme is positive beyond solver precision, so that a certificate is to be expected, and \emph{undecided} if this optimum is $0$ to solver precision, as in Remark~\ref{rem:limits}.
Here the programme maximises the margin of a certificate of the given form under a normalisation, the margin being the sum of the slacks in (i) and (ii) of Lemma~\ref{lem:flagcert}.
As the margin scales with the certificate, a certificate of this form exists exactly when this optimum is positive.
So an optimum of $0$ to solver precision, or a slightly negative one as in Remark~\ref{rem:limits}, means that, up to solver precision, none exists.
\begin{itemize}
\item Without the localised squares, the same method gives certificates for $n=218$, $212$, $208$ and $204$, of which only the last is in the ancillary files, and in floating point its relaxation is decided down to $n=200$.
\item The five-vertex model of Theorem~\ref{thm:r34five} (the \emph{mixture model}, described under ``Certificates and margins'' below), in which the profile of $x_0$ is an unknown mixture of admissible profiles, also gives $\R(4,4,4)\le 200$.
Its certificate, \texttt{R444\_200.json}, is in the folder \texttt{flag5\_r34k/} of the ancillary files.
\item The four-vertex method of Section~\ref{sec:444} with several squares reaches $219$, and its full semidefinite relaxation is undecided at $218$.
\end{itemize}

\paragraph{Why it stops here.}
In the certificate the multiplier weight sits on the cap rows
\[
|N_i(u)\cap N_i(v)|\le 17
\]
for an edge $uv$ of colour $i$, which come from $\R(4,4)=18$.
The degree cap $76$ carries almost no weight.
Correspondingly, a better bound for $\R(3,4,4)$ alone did not help in our runs.
With the hypothesis $\R(3,4,4)\le 75$ in place of $77$ the relaxation at $n=197$ is still undecided (we did not rerun this model with our bound $\R(3,4,4)\le 73$ of Theorem~\ref{thm:r34five}), and the mixture model stops at $200$ with $\R(3,4,4)\le 73$ as with $77$.

At $n=197$ the relaxation is undecided to solver precision, and it stays so with either of two six-vertex extensions, localised squares at triangles and two-vertex flags at an edge.
Even under the extra hypothesis that every vertex has a regular profile, with entries differing by at most $1$, the relaxation is decided numerically at $197$ but not at $196$.
So conditioning on profiles could gain at most one order.

\paragraph{Certificates and margins.}
For the certificates of Theorem~\ref{thm:r34five}, and the mixture-model certificate for $\R(4,4,4)\le 200$ above, the star term has a coefficient for each of the $3^4$ colour patterns of the star, and the search treats the law of the profile of $x_0$ as an unknown mixture of the admissible profiles.
Each certificate is a function $F$, the star term plus nonnegative multiples of cap and lower rows minus nonnegative multiples of squares, with
\[
\min\mathrm{Sym}\,F>\max J_G,
\]
the minimum over the colourings of $K_5$ without a forbidden clique and the maximum over the admissible profiles.
For any constant $C$ between these two values, $C-F$, with $C$ absorbed into its star term, satisfies (i) and (ii) of Lemma~\ref{lem:flagcert}.
The \emph{margin} is
\[
\min\mathrm{Sym}\,F-\max J_G,
\]
the sum of the slacks in (i) and (ii).

\begin{theorem}[certificates checked by computer]\label{thm:r34five}
\begin{gather*}
\R(3,4,4)\le 73,\qquad \R(3,4,5)\le 137,\qquad \R(3,5,5)\le 316,\\
\R(3,4,6)\le 247,\qquad \R(3,4,7)\le 409,\\
\R(3,5,6)\le 622\qquad\text{and}\qquad \R(3,4,8)\le 626.
\end{gather*}
With the starred inputs of Table~\ref{tab:inputs},
\begin{gather*}
\R(3,4,6)\le 245,\qquad \R(3,4,7)\le 406,\\
\R(3,5,6)\le 619\qquad\text{and}\qquad \R(3,4,8)\le 625.
\end{gather*}
\end{theorem}

\begin{proof}
One certificate for each bound, with $s=5$, checked in exact arithmetic over all $3^{10}$ labelled colourings of $K_5$:
\begin{center}\small
\begin{tabular}{lrrrrl}
\toprule
target & $n$ & orbits & profiles & margin & new bounds used\\
\midrule
$(3,4,4)$ & $73$ & $273$ & $10$ & $2.2\cdot10^{-5}$ & none\\
$(3,4,5)$ & $137$ & $551$ & $231$ & $5.4\cdot10^{-7}$ & none\\
$(3,5,5)$ & $316$ & $287$ & $6$ & $1.5\cdot10^{-7}$ & $73$, $137$\\
$(3,4,6)$ & $247$ & $552$ & $171$ & $4.1\cdot10^{-7}$ & $73$, $137$\\
$(3,4,7)$ & $409$ & $552$ & $378$ & $5.5\cdot10^{-8}$ & $73$, $137$, $247$\\
$(3,5,6)$ & $622$ & $566$ & $378$ & $3.6\cdot10^{-8}$ & $73$, $137$, $247$, $316$\\
$(3,4,8)$ & $626$ & $552$ & $1128$ & $2.7\cdot10^{-8}$ & $73$, $137$, $247$, $409$\\
\bottomrule
\end{tabular}
\end{center}
The table rows are for the bounds with published inputs.
\begin{itemize}
\item $\R(3,4,4)\le 73$ uses only $\R(3,3)=6$, $\R(3,4)=9$, $\R(4,4)=18$ and $\R(3,3,4)=30$, and needs the exchange of blue and green, which lies in $\Gamma$.
\item $\R(3,4,5)\le 137$ uses in addition $\R(3,5)=14$, $\R(4,5)=25$, $\R(3,3,3)=17$, $\R(3,3,5)\le 57$ and $\R(3,4,4)\le 77$.
\item The others also use
\begin{itemize}
\item the values $\R(3,k)$ with $k\le 8$,
\item the published bounds $\R(5,5)\le 46$, $\R(4,6)\le 41$, $\R(4,7)\le 61$, $\R(4,8)\le 84$ and $\R(5,6)\le 87$,
\item the bounds of the theorem listed in the table, and
\item where they enter a cap, the hand bounds of Theorems~\ref{thm:336}, \ref{thm:337} and~\ref{thm:338}.
\end{itemize}
\end{itemize}
With $\R(3,4,4)\le 77$ in place of $73$, the $(3,5,5)$ certificate remains valid, with the same margin, once its one cap constant that involves $\R(3,4,4)$ is recomputed, from $72$ to $76$.
The stored certificate has the constant $72$; we checked that, with $76$ in its place, the same checker accepts it without the input $\R(3,4,4)\le 73$, but this variant is not among the ancillary files.
The $(3,4,8)$ certificate uses $409$ only through $\R(3,4,7)\le 411$, so that $D_3=410$.

The variants with starred inputs replace the published two-colour bounds by those of Table~\ref{tab:inputs}, and the one for $\R(3,4,8)$ also uses $\R(3,4,6)\le 245$.
\end{proof}

For comparison, $\R(3,4,5)\le 137$ and \eqref{eq:recurrence} give only $\R(3,5,5)\le 46+137+137-1=319$.
Similarly, $\R(3,4,4)\le 73$ gives only $\R(4,4,4)\le 3\cdot 73-1=218$.
Theorem~\ref{thm:r34five} supersedes Theorems~\ref{thm:345}, \ref{thm:346} and~\ref{thm:355} and Proposition~\ref{prop:supp}, whose hand proofs remain as short independent arguments for weaker bounds.

For $\bt=(3,3,k)$, Lemma~\ref{lem:flagcert} gives both exclusions and valid lower bounds $m_k(n)$, and the procedure of Section~\ref{sec:ladder} can take the larger of these and its own at every order.

\begin{theorem}[certificates checked by computer]\label{thm:flag678}
$\R(3,3,6)\le 86$, $\R(3,3,7)\le 125$ and $\R(3,3,8)\le 174$.
\end{theorem}

\begin{proof}
Three certificates with $s=5$, for $(k,n)=(6,86)$, $(7,125)$ and $(8,174)$, checked in this order.
They have $64$, $37$ and $71$ rows, all cap and lower rows except one upper floor row in the third.
They have $121$, $144$ and $106$ squares, mostly at triangles.
None has a profile row of the recursive kind.
So they use only the inputs of Table~\ref{tab:inputs} and, at $k=7$ and $8$, the bounds just proved, through the caps.

They are checked in exact arithmetic over the $193$ orbits of $S_5\times\mathbb Z_2$ on the $29{,}197$ labelled colourings of $K_5$ with no red and no blue triangle, and over the admissible profiles.
The least values in (ii) are about $1.9\cdot10^{-7}$, $3.1\cdot10^{-8}$ and $1.5\cdot10^{-8}$.
\end{proof}

Theorem~\ref{thm:flagladder} lowers the third bound to $173$.
So Theorems~\ref{thm:flag678} and~\ref{thm:flagladder} supersede Theorems~\ref{thm:336}, \ref{thm:337} and~\ref{thm:338}, whose hand proofs remain as short independent arguments for weaker bounds.
The certificates of Theorem~\ref{thm:r34five} use these weaker bounds as caps.

\begin{theorem}[computer-assisted]\label{thm:flagladder}
$\R(3,3,k)\le 173$, $241$, $309$, $401$, $505$, $620$, $748$ and $892$ for $k=8,\dots,15$.
\end{theorem}

\begin{proof}[Verification]
A single chain, the procedure of Section~\ref{sec:ladder} with valid lower bounds also taken from certificates, is checked level by level by a program that does not import the generating code.
At each level $k=3,\dots,15$ the program works as follows.
\begin{enumerate}
\item It first checks the certificates of that level.
A recursive profile row is accepted only with values already accepted at level $k-1$, or with the bound $\lceil z\max(0,z-1-C_{33})/2\rceil$ at green degree $z$, as each vertex of $N_3(v)$ has at most $C_{33}$ green neighbours in $N_3(v)$.
\item It then checks the rows of the ladder at level $k$.
Each claimed $m_k(n)$ must be at most the largest of
\begin{itemize}
\item $\lceil nh_{\min}/2\rceil$,
\item the bounds from the recorded supporting lines of the functions $F$ and $\Psi$ of Lemmas~\ref{lem:commonnbr} and~\ref{lem:density},
\item Lemma~\ref{lem:delete}(a) and~(b), and
\item an accepted certificate value at $(k,n)$.
\end{itemize}
\item Finally it checks the exclusion of $B_k$, by a certificate or by $\Psi>0$ at every admissible $h$.
\end{enumerate}
No bound on $\R(3,3,k)$ with $k\ge 6$ is an input: each $B_k$ is treated as infinite until the chain excludes it, so every bound $B_j$ with $j\ge 6$ used at a level is proved earlier in the chain.

The chain has
\begin{itemize}
\item $116$ \emph{density certificates}, of the second form in Lemma~\ref{lem:flagcert}, which give valid lower bounds $m_k(n)$, with $s=4$ at levels $3$ to $6$,
\item an exclusion of $87$ with $s=4$, and
\item $38$ certificates with $s=5$: densities at levels $6$, $7$ and $8$, the exclusions of Theorem~\ref{thm:flag678}, and a few that add nothing.
\end{itemize}
For $8\le k\le 15$ the final exclusion is by $\Psi$.
Without certificates, the same checker reproduces Table~\ref{tab:ladder} exactly.
\end{proof}

\begin{table}[htbp]
\centering\small
\begin{tabular}{l*{10}{c}}
\toprule
$k$ & 6 & 7 & 8 & 9 & 10 & 11 & 12 & 13 & 14 & 15\\
\midrule
Sections~\ref{sec:three} and~\ref{sec:ladder} & 88 & 129 & 179 & 247 & 315 & 407 & 512 & 628 & 756 & 901\\
with certificates & 86 & 125 & 173 & 241 & 309 & 401 & 505 & 620 & 748 & 892\\
\bottomrule
\end{tabular}
\caption{Upper bounds on $\R(3,3,k)$.}
\label{tab:flagladder}
\end{table}

\begin{remark}[Limits]\label{rem:limits}
At one order below each bound of Theorem~\ref{thm:r34five}, except the starred $245$, below which $244$ was not tried, the optimum of the semidefinite programme is $0$ to solver precision, between $-8\cdot10^{-9}$ and $-6\cdot10^{-10}$.
In particular we could not exclude $(3,4,4;72)$.
The optimal profile mixture puts weight $0.89$ on $(17,27,27)$, and adding the exact pattern counts of the unique $(4,4;17)$-colouring \cite{Rad26} inside $N_1(v)$ when $d_1(v)=17$ still gives no certificate.

For $\R(3,3,k)$ the five-vertex cutting-plane searches at $(3,3,6;85)$, $(3,3,7;124)$ and $(3,3,8;172)$ were stopped without a certificate.
When they were stopped, the cutting-plane relaxations still had points satisfying every cut with slack about $3\cdot10^{-7}$, $6.5\cdot10^{-8}$ and $2.5\cdot10^{-8}$; a certificate appears only once this slack becomes negative.
So these searches are inconclusive: they found no certificate, but did not show that none exists.
So these searches are inconclusive: they found no certificate, but did not show that the optimum is $0$.
Neither the four- nor the five-vertex search excludes $(3,3,5;56)$.
With $s=4$ nothing is gained for $\R(3,3,k)$ above $k=6$.

The search of Theorem~\ref{thm:r34five} gives no certificate at orders where colourings are known, for instance $(3,3,4;29)$, $(3,4,4;54)$ and $(4,4,4;127)$.
\end{remark}

\section{Lower bounds from split circulants}\label{sec:lower}

The witnesses for $\R(3,3,k)$ with $k\le 14$ in Table~\ref{tab:lower} are circulant colourings.
The vertex set is $\mathbb Z_n$, and the \emph{distance} of a pair $\{u,v\}$ is the $d\in\{1,\dots,\lfloor n/2\rfloor\}$ with $u-v\equiv\pm d\pmod n$.
Given disjoint sets $\Delta_1,\Delta_2\subseteq\{1,\dots,\lfloor n/2\rfloor\}$, colour $\{u,v\}$ red if its distance lies in $\Delta_1$, blue if it lies in $\Delta_2$, and green otherwise.
Thus red is the Cayley graph of $\mathbb Z_n$ with connection set
\[
\pm\Delta_1:=\{\pm a\bmod n:a\in\Delta_1\}.
\]

As translations preserve colours, red has a triangle if and only if it has one through $0$, say $\{0,x,x+y\}$.
That is, it has one if and only if $x,y,x+y\in\pm\Delta_1$ for some $x,y$ (the three vertices are then distinct, as $0\notin\pm\Delta_1$).
So red is triangle-free if and only if $\pm\Delta_1$ is \emph{sum-free}: $x+y\notin\pm\Delta_1$ for all $x,y\in\pm\Delta_1$, including $x=y$.
Likewise for blue and $\pm\Delta_2$.
A green $K_k$ is excluded separately, by clique search.

\paragraph{Colour splitting.}
The colourings were found by \emph{colour splitting}, the method of Exoo and Tatarevic \cite[\S3.5]{ET15}.
It starts from a two-colour circulant on $n$ vertices, the \emph{seed}, whose second colour is $K_k$-free.
A SAT solver, CaDiCaL~1.5 through PySAT \cite{IMM18}, searches for a partition of the distances of its first colour into classes $\Delta_1$ and $\Delta_2$ with $\pm\Delta_1$ and $\pm\Delta_2$ sum-free.
Green is then the seed's second colour, so the split is a $(3,3,k;n)$-colouring.
As $\R(3,3)=6$, a split can exist only if the first colour is $K_6$-free.

\paragraph{The seeds.}
Our seeds are $(5,k)$ circulants of Kuznetsov, whose first colour is even $K_5$-free, though a split still need not exist.
For $k=9,11,12,14$ they are those of \cite[Table~2]{Kuz16}.
For $k=13$ the circulant there, on $232$ vertices, admits no split, even without circulant symmetry: with one variable per edge the condition is unsatisfiable, already on any $51$ consecutive vertices.
We use instead a $231$-vertex circulant $(5,13)$-colouring from the program distributed with the source files of \cite{Kuz16}.

We chose these seeds because their orders exceed those of the witnesses for the survey's lower bounds on $\R(3,3,k)$.
For $k=9,11,12,14$ the new bound $\R(3,3,k)\ge n+1$ equals the survey's lower bound on $\R(5,k)$.
For $k=13$ it is one less than $\R(5,13)\ge 233$.

For $k=10$ no known $(5,10)$-colouring is as large as the survey's $149$-vertex witness for $\R(3,3,10)$.
The largest, of Exoo and Tatarevic \cite{ET15} (see \cite[Table~3]{Kuz16}), has $148$ vertices and is not a circulant.
Although it splits edge by edge, it gives only $149$.
For $k=15$ we instead add vertices to another split of the seed for $k=14$ (Section~\ref{sec:extend}).

The splits are ours.
The verification checks all three colours directly, so the validity of the witnesses does not depend on the seed or the solver.

\begin{theorem}[witnesses checked by computer]\label{thm:lower}
Let $(n,k)$ be one of the pairs $(132,9)$, $(182,11)$, $(202,12)$, $(231,13)$ and $(266,14)$.
Colour the edges of the complete graph on $\mathbb Z_n$ as above, with $\Delta_1$ and $\Delta_2$ the distance sets listed below for this pair.
This colouring has no red or blue triangle and no green $K_k$.
Hence $\R(3,3,k)\ge n+1$ (Table~\ref{tab:lower}).
\end{theorem}

Each item gives $|\Delta_1|$, $|\Delta_2|$ and, as a check on transcription, the degree profile $(d_1,d_2,d_3)$, which is common to all vertices.
Each vertex has two neighbours at every distance $d<n/2$ and, when $n$ is even, one at distance $n/2$, as $n/2\equiv-n/2\pmod n$.
So a colour with distance set $\Delta$ has degree $2|\Delta|$, less $1$ if $n$ is even and $n/2\in\Delta$.
Also $d_3=n-1-d_1-d_2$.
\begin{itemize}[leftmargin=1.2em]\small
\item $(n,k)=(132,9)$: $\Delta_1=\{$5, 9, 11, 12, 15, 19, 29, 32, 36, 39, 46, 53, 56, 59, 63, 66$\}$, $\Delta_2=\{$6, 16, 17, 20, 27, 28, 30, 31, 38, 41, 42$\}$; $|\Delta_1|=16$, $|\Delta_2|=11$, $n/2=66\in\Delta_1$, profile $(31,22,78)$.
\item $(n,k)=(182,11)$: $\Delta_1=\{$29, 36, 43, 50, 54, 61, 63, 68, 70, 75, 77, 80, 81, 84, 87, 88, 91$\}$, $\Delta_2=\{$4, 5, 7, 13, 21, 22, 30, 31, 39, 40, 48, 64, 65, 66, 74, 82, 83$\}$; $|\Delta_1|=17$, $|\Delta_2|=17$, $n/2=91\in\Delta_1$, profile $(33,34,114)$.
\item $(n,k)=(202,12)$: $\Delta_1=\{$2, 9, 15, 16, 22, 23, 28, 35, 41, 48, 49, 52, 60, 66, 73, 78, 79, 85, 86, 91, 92, 99$\}$, $\Delta_2=\{$1, 7, 10, 19, 27, 33, 36, 45, 51, 59, 62, 71, 76, 88, 97, 101$\}$; $|\Delta_1|=22$, $|\Delta_2|=16$, $n/2=101\in\Delta_2$, profile $(44,31,126)$.
\item $(n,k)=(231,13)$: $\Delta_1=\{$3, 4, 5, 12, 13, 20, 21, 29, 31, 37, 46, 48, 54, 65, 76, 82, 93, 104, 110$\}$, $\Delta_2=\{$2, 6, 10, 17, 18, 39, 42, 43, 47, 51, 55, 63, 67, 71, 75, 79, 83, 87, 91, 95, 99, 103$\}$; $|\Delta_1|=19$, $|\Delta_2|=22$, $n$ odd, profile $(38,44,148)$.
\item $(n,k)=(266,14)$: $\Delta_1=\{$6, 8, 9, 10, 11, 13, 39, 41, 42, 43, 44, 64, 65, 67, 68, 69, 89, 90, 92, 93, 94, 122, 123, 124, 125, 127$\}$, $\Delta_2=\{$1, 46, 49, 52, 55, 58, 61, 66, 96, 99, 102, 105, 108, 111, 114, 117, 120, 126, 129, 131, 133$\}$; $|\Delta_1|=26$, $|\Delta_2|=21$, $n/2=133\in\Delta_2$, profile $(52,41,172)$.
\end{itemize}

\begin{proof}[Verification]
\emph{The sum-free check.}
Red and blue are triangle-free by the sum-free criterion above, a finite check on $\pm\Delta_1$ and $\pm\Delta_2$.

\emph{Green cliques.}
As translations preserve colours, every green clique has a translate containing $0$, so it suffices to find the largest green clique through $0$.
An exact branch-and-bound bitset search for cliques through vertex~$0$, \texttt{clique\_check.cpp}, gives green clique numbers $8$, $10$, $11$, $12$ and $13$ for $k=9$, $11$, $12$, $13$ and $14$, that is, $k-1$.
So there is no green $K_k$.

\emph{An independent check.}
A second exact search, \texttt{check\_colouring\_full.cpp}, which shares no code with \texttt{clique\_check.cpp}, works on the full adjacency matrices without using translation invariance.
It confirms that there is no green $K_k$ and no red or blue triangle.
The bounds do not need the green $K_{k-1}$, but it shows that the colouring is not $(3,3,k-1)$-avoiding, so it does not also give $\R(3,3,k-1)\ge n+1$.
The colourings and programs, including the sum-free check, are in the ancillary files (Section~\ref{sec:verify}).
\end{proof}

\subsection{Extending a split circulant}\label{sec:extend}

A witness for $k+1$ can be built from one for $k$ by adding vertices.
Let $H$ be a colouring of $K_n$ with no red or blue triangle and with green clique number $q$, and let $F$ be a colouring of $K_{n+r}$ whose restriction to the first $n$ vertices is $H$.
We call the other $r$ vertices \emph{new}.

For a new vertex $x$ let $A_x$ be the set of vertices of $H$ joined to $x$ in red or blue.
For a set $S$ of new vertices let
\[
A_S=\bigcup_{x\in S}A_x.
\]
A green clique of $F$ is $S\cup T$, where $S$ is a green clique of new vertices and $T$ a green clique of $H-A_S$.
A vertex $w$ of $H$ \emph{certifies} $S$ if $w$ and all its red and blue neighbours in $H$ lie in $A_S$.

\begin{lemma}\label{lem:cert}
Let $W$ be a set of certifiers of $S$ that is a green clique of $H$.
Then every green clique of $H-A_S$ has at most $q-|W|$ vertices.
\end{lemma}

\begin{proof}
Let $T$ be such a clique.
It is disjoint from $W$, as $W\subseteq A_S$.
A vertex of $T$ is outside $A_S$, so it is not a red or blue neighbour of any $w\in W$, and is therefore a green neighbour of every $w\in W$.
So $T\cup W$ is a green clique of $H$, and $|T|+|W|\le q$.
\end{proof}

So $F$ has green clique number at most $q+1$ if every green clique $S$ of new vertices satisfies
\[
|S|\le 1+\omega(W_S),
\]
where $W_S$ is the set of certifiers of $S$ and $\omega(W_S)$ is the green clique number of $H$ on $W_S$.
Checking this needs no clique search in $F$.
It runs over the green cliques of new vertices, and needs only the green clique numbers of the sets $W_S$.

The simplest new vertex is a \emph{copy}.
The copy $y_v$ of a vertex $v$ of $H$ has the colours of $v$ towards $H-v$, and $vy_v$ is green.
So $A_{y_v}$ is the set of red and blue neighbours of $v$.
Then $v$ certifies $\{x,y_v\}$ for every new vertex $x$ with $v\in A_x$, and $\{y_u,y_v\}$ whenever $uv$ is red or blue.

\begin{theorem}[witness checked by computer]\label{thm:ext15}
$\R(3,3,15)\ge 278$.
\end{theorem}

\begin{proof}
\emph{Step 1: the colouring $H$.}
Let $H$ be the circulant colouring of $\mathbb Z_{266}$ with $\Delta_1=\{$1, 42, 46, 49, 52, 55, 58, 61, 64, 96, 99, 102, 105, 108, 111, 114, 117, 120, 127, 129, 131, 133$\}$ and $\Delta_2=\{$6, 8, 9, 10, 11, 13, 39, 41, 43, 44, 65, 66, 67, 68, 69, 89, 90, 92, 93, 94, 122, 123, 124, 125, 126$\}$; $|\Delta_1|=22$, $|\Delta_2|=25$, $n/2=133\in\Delta_1$, profile $(43,50,172)$.
Both $\pm\Delta_1$ and $\pm\Delta_2$ are sum-free, and $\Delta_1\cup\Delta_2$ is the same set as for the $(266,14)$ witness of Theorem~\ref{thm:lower}.
So $H$ is another split of the same seed, with the same green graph, and $q=13$.

\emph{Step 2: the new vertices.}
Add eleven new vertices:
\begin{itemize}
\item the copies $y_0,y_5,y_{70},y_{131},y_{140},y_{201}$ of the \emph{anchors} $0,5,70,131,140,201$; and
\item five vertices $x_1,\dots,x_5$, each green to every vertex of $H$ except as follows:
\begin{itemize}
\item $x_1$ is red to $0$ and $70$ and blue to $201$;
\item $x_2$ is red to $131$ and $201$ and blue to $70$;
\item $x_3$ is red to $5$ and blue to $140$;
\item $x_4$ is red to $140$ and blue to $201$;
\item $x_5$ is red to $5$ and blue to $70$.
\end{itemize}
\end{itemize}

\noindent\begin{minipage}{\linewidth}
\emph{Step 3: colours among the new vertices.}
Colour the pairs of new vertices, taken in the order $y_0,y_5,y_{70},y_{131},y_{140},y_{201},x_1,\dots,x_5$, by the following rows.
Row $i$ lists the colours from the $i$th new vertex to the eleven new vertices in this order (r red, b blue, g green; a dot marks the diagonal):
\begin{center}\small\ttfamily
.rbggbgbrrb\quad r.bgrbbrgbg\quad bb.rbrggbrg\quad ggr.bbbgrbr\quad grbb.rrbggb\quad bbrbr.ggbgr\\
gbgbrg.brrb\quad brggbgb.rbr\quad rgbrgbrr.bb\quad rbrbggrbb.r\quad bggrbrbrbr.
\end{center}
\end{minipage}

So $x_1,\dots,x_5$ span a red pentagon $x_1x_3x_2x_5x_4$ and a blue pentagram, and $x_i$ is green to $y_v$ exactly when $x_iv$ is red or blue.

\emph{Step 4: no red or blue triangle.}
The resulting colouring $F$ of $K_{277}$ has no red or blue triangle, by a direct check.

\emph{Step 5: green cliques.}
The green edges among the new vertices are $y_0y_{131}$, $y_0y_{140}$, $y_5y_{131}$ and the twelve pairs $x_iy_v$ with $x_iv$ red or blue.
They span no green triangle.
Each has a certifier:
\begin{itemize}
\item $0$ for the first two and $5$ for the third, by the rule for copies, as the pairs $\{0,131\}$, $\{0,140\}$ and $\{5,131\}$, at distances $131$, $126$ and $126$, are red or blue in $H$; and
\item $v$ for $x_iy_v$.
\end{itemize}
Let $S\cup T$ be a green clique of $F$, as before Lemma~\ref{lem:cert}, so that $|S|\le 2$.
If $|S|\le 1$, then $|S\cup T|\le 1+q$.
If $|S|=2$, then $S$ has a certifier $w$, and Lemma~\ref{lem:cert} with $W=\{w\}$ gives $|T|\le q-1$.
Either way $|S\cup T|\le q+1=14$.
So $F$ is a $(3,3,15;277)$-colouring.
\end{proof}

\begin{proof}[Verification]
\emph{The matrix and the certifiers.}
One program rebuilds the $277\times 277$ matrix from $\Delta_1$, $\Delta_2$, the anchors and the rows above, so its first $266$ vertices carry $H$.
Another, \texttt{certcheck.py}, reads the matrix, checks that red and blue have no triangle, and for every green clique $S$ of new vertices computes the certifiers and checks the condition after Lemma~\ref{lem:cert}.
With $q=13$ from Theorem~\ref{thm:lower}, this is a complete check.

\emph{Independent searches.}
Independently, and without the lemma, \texttt{check\_colouring\_full.cpp} of Theorem~\ref{thm:lower} and a second exact branch-and-bound search on the full matrix find no red or blue triangle and no green $K_{15}$.
The second search, and a third that computes $\max_S\big(|S|+\omega(H-A_S)\big)$ exactly over the decomposition $S\cup T$, give green clique number $14$.
\end{proof}

\paragraph{The search and its limits.}
The anchors and the colours among the new vertices were found by a SAT search over copies of a set of anchors together with vertices that are green to $H$ outside the anchors.
In none of the $416$ circulant splits of the seed, counted up to exchanging red and blue, does this family give more than eleven new vertices.

Call vertices $u,v$ of $H$ \emph{compatible} if they have no common red neighbour or no common blue neighbour, as holds whenever $uv$ is red or blue.
Two copies $y_u,y_v$ with $uv$ green cannot be green to each other, as the common green neighbourhood of $u$ and $v$ contains a green $K_{11}$.
They can be red only if $u$ and $v$ have no common red neighbour, and blue only if they have no common blue neighbour.
So the anchors are pairwise compatible.
In our enumeration of the splits, at most six vertices of $H$ are pairwise compatible, so there are at most six anchors and six copies.

Two new vertices of the second kind cannot be green to each other either, as $H$ minus the at most six anchors still contains a green $K_{13}$.
Indeed, as $H$ is vertex-transitive, a set meeting every green $K_{13}$ has at least $266/13>20$ vertices.
So every pair of them is red or blue, and as $\R(3,3)=6$ there are at most five of them.

\section{Verification, literature and limitations}\label{sec:verify}

\paragraph{Ancillary files.}
The ancillary files of this submission, also in the directory \texttt{ancillary/} of the repository \url{https://github.com/Lachy-Dauth/ramsey-neighbourhood-counting}, contain programs that check every computer-assisted result of the paper.
Each certificate is checked by a program written from the definitions, which does not import the code that found it.
The Python, C and C++ programs use only their standard libraries, except that the six-vertex checker in \texttt{flag5\_r444/} and its negative controls use NumPy for integer counts, and the C programs of \texttt{two\_colour\_r4k/}, for the companion paper, also use POSIX threads.
Statements about searches and negative experiments (for instance in Sections~\ref{sec:five} and~\ref{sec:lower} and in Remark~\ref{rem:limits}) report our runs and are not certified by these files.

The command \texttt{sh run\_all.sh} runs the quick checks in a few minutes, \texttt{sh run\_all.sh --full} adds the long checks, and \texttt{sh run\_all.sh --controls} the negative and positive controls.
The checks of Theorems~\ref{thm:r444five}, \ref{thm:r34five} and~\ref{thm:flagladder}, the second checker for Theorem~\ref{thm:r444five}, the full-matrix searches of \texttt{check\_colouring\_full.cpp} and the direct clique searches for Theorem~\ref{thm:ext15} run only with \texttt{--full}.
Of Section~\ref{sec:five}, the quick mode checks Theorem~\ref{thm:flag678} and the weaker bound $\R(4,4,4)\le 204$.
\begin{itemize}[leftmargin=1.2em]\raggedright
\item \texttt{verify\_hand\_proofs.py}: the main steps of the proofs in Sections~\ref{sec:three} and~\ref{sec:types} (some values quoted only in remarks are not checked);
\item \texttt{bootstrap.py}, which runs the procedure of Section~\ref{sec:ladder} and writes its certificate \texttt{tables.json}, and \texttt{verify\_bootstrap.py}, a checker that shares no code with it: Theorem~\ref{thm:ladder};
\item \texttt{verify\_cascade.py} and \texttt{verify\_cascade\_published\_only.py}, with and without the unpublished inputs: the program in the proof of Proposition~\ref{prop:supp};
\item \texttt{verify\_r444.py}: Lemma~\ref{lem:444id}, Theorem~\ref{thm:444a}, and identity \eqref{eq:id225} with the values of $J$ in Theorem~\ref{thm:444b};
\item \texttt{verify\_r444\_squares.py}: the four-vertex certificates with several squares for $n=219,\dots,224$;
\item \texttt{flag5\_r444/}: Theorem~\ref{thm:r444five}, with a second checker over six-vertex orbits;
\item \texttt{flag5\_r34k/}: Theorem~\ref{thm:r34five} and the mixture-model certificate for $\R(4,4,4)\le 200$, with a manifest that checks which bounds each certificate uses;
\item \texttt{flag\_r33k/}: Theorems~\ref{thm:flag678} and~\ref{thm:flagladder}, the whole chain of certificates and ladder rows;
\item \texttt{circulants.json}, \texttt{check\_circulants.py}, \texttt{clique\_check.cpp} and \texttt{check\_colouring\_full.cpp}: Theorem~\ref{thm:lower}; \texttt{check\_circulants.py} regenerates the adjacency matrices of the five witnesses from $\Delta_1$ and $\Delta_2$ (the lists \texttt{A} and \texttt{B} of \texttt{circulants.json}) and prints the first $16$ hexadecimal digits of their SHA-256 hashes;
\item \texttt{extension\_r3315/}: Theorem~\ref{thm:ext15}, with \texttt{certcheck.py} for the check through Lemma~\ref{lem:cert}.
\end{itemize}

\paragraph{Tests and negative controls.}
All these checks were run, together with negative controls (perturbed certificates, constants changed by one, the order changed by one), all of which were rejected.
For the flag certificates of Section~\ref{sec:five}, the rows were also evaluated on explicit colourings avoiding the target, where every row holds, and the square expansions were compared exactly with the squares on explicit colourings.
These tests check the implementation.
At the orders of the theorems no colouring exists, and the validity of the rows rests on their proofs in Section~\ref{sec:cert}.

\paragraph{Formalisation in Lean.}
Two Lean~4 files in \texttt{lean/}, \texttt{LocalCombinatorics.lean} and \texttt{ConventionalDefinitions.lean}, which need Mathlib, formalise with no unproved steps (no \texttt{sorry}) two facts used in Lemma~\ref{lem:apex}, on which all our hand bounds for $\R(3,3,k)$ rest.
\begin{itemize}
\item The first is the cap $|N_1(v)\cap N_2(u)|\le k-1$ for $u\in N_1(v)$ in a $(3,3,k)$-colouring, and its sum over $u\in N_1(v)$, whose left side is $2e_2(N_1(v))$.
So $2e_2(N_1(v))\le(k-1)d_1(v)$ (Lemma~\ref{lem:caps} and Corollary~\ref{cor:handshake}).
\item The second is the pointwise identity behind $T(H)=\sum_v\big(e_2(N_1(v))+e_1(N_2(v))\big)$ (Lemma~\ref{lem:corners}).
In a colouring with no red or blue triangle, each triangle of $H$ contributes exactly $1$ to the sum, at its apex, and every other triangle contributes $0$.
\end{itemize}
Nothing else is formalised.

\paragraph{Literature.}\label{sec:lit}
Comparison values are from revision~18 of the survey \cite{Rad26} and from \cite{Boz26}.
We searched arXiv and the web up to 23 October 2026 for earlier statements of our new bounds, and found none.

The search included the flag-algebra tables of Lidick\'y and Pfender \cite{LP21}.
Among three-colour numbers with every $t_i\ge 3$, these bound only $\R(3,3,4)$, $\R(3,3,5)$, $\R(3,4,4)$ and $\R(3,4,5)$, by $32$, $57$, $85$ and $406$.
None is below the bound known at the time.
For $\R(3,3,5)$ they match $57$.

Such a search cannot establish priority, and we would welcome pointers to any earlier source.

\paragraph{Limitations.}
The procedure of Section~\ref{sec:ladder} alone improves neither $\R(3,3,5)\le 57$ nor the hand bounds $88$, $129$ and $179$ (Remark~\ref{rem:sanity}).
With the certificates of Section~\ref{sec:five} it gives Table~\ref{tab:flagladder}.
No method here improves $\R(3,3,5)\le 57$, on which every bound for $\R(3,3,k)$ with $k\ge 6$ depends.

One order below the bounds of Theorems~\ref{thm:r444five} and~\ref{thm:r34five}, except the starred $245$ (the order $244$ was not tried), the five-vertex relaxations have optimum $0$ to solver precision.
We also found no five-vertex certificate for $(3,3,6;85)$, $(3,3,7;124)$ or $(3,3,8;172)$ (Section~\ref{sec:five} and Remark~\ref{rem:limits}).
Further progress for these numbers needs larger flags, new valid inequalities, or better inputs.

Searches in restricted families, namely circulant colourings and block and extension constructions, found no $(3,3,7;85)$-colouring and no $(3,4,4;55)$-colouring, which would improve the lower bounds $\R(3,3,7)\ge 85$ and $\R(3,4,4)\ge 55$ listed in the survey, respectively.
A failure in a restricted family proves nothing in general.
\section*{Acknowledgements}
\pdfbookmark[1]{Acknowledgements}{acknowledgements}
Much of the work reported here was carried out with two AI systems, used under the direction of the author: OpenAI's Codex, based on GPT-6, in the first stage of the project, and Claude Opus~5.5, a large language model developed by Anthropic, in the later stages.
Between them they searched for the arguments and certificates, wrote the search programs, the checking programs and the Lean files, and drafted the text.
Of the search programs only \texttt{bootstrap.py} is in the ancillary files; all the checking programs are.
Every certificate and witness is checked in exact arithmetic by the programs in the ancillary files, which can be run independently of the model.
The author takes full responsibility for the contents of this paper.

\pdfbookmark[1]{References}{references}
\begin{thebibliography}{99}\small
\bibitem{AR05} N.~Alon and V.~R\"odl, Sharp bounds for some multicolor Ramsey numbers, \emph{Combinatorica} 25 (2005), 125--141.
\bibitem{ACM26} B.~Andrades, M.~Campos and R.~Morris, On the multicolour Ramsey numbers $R(3,3,k)$, arXiv:2609.28455, 2026.
\bibitem{Ang25} V.~Angeltveit, $R(3,10)\le 41$, \emph{Electron.\ J.\ Combin.} 32(4) (2025), \#P4.30.
\bibitem{AM26} V.~Angeltveit and B.~D.~McKay, $R(5,5)\le 46$, \emph{J.\ Graph Theory} 112 (2026), no.~3, 198--208; arXiv:2409.15709.
\bibitem{Boz26} L.~Boza, New upper bounds for the classical Ramsey numbers $R(4,4,4)$, $R(3,4,5)$ and $R(3,3,6)$, arXiv:2603.10851, 2026.
\bibitem{CFIM16} M.~Codish, M.~Frank, A.~Itzhakov and A.~Miller, Computing the Ramsey number $R(4,3,3)$ using abstraction and symmetry breaking, \emph{Constraints} 21 (2016), 375--393.
\bibitem{DC26} L.~Dauth, New bounds for small two- and four-colour Ramsey numbers, companion paper, 2026.
\bibitem{ET15} G.~Exoo and M.~Tatarevic, New lower bounds for 28 classical Ramsey numbers, \emph{Electron.\ J.\ Combin.} 22(3) (2015), \#P3.11.
\bibitem{GR13} J.~Goedgebeur and S.~P.~Radziszowski, New computational upper bounds for Ramsey numbers $R(3,k)$, \emph{Electron.\ J.\ Combin.} 20(1) (2013), \#P30.
\bibitem{Goo59} A.~W.~Goodman, On sets of acquaintances and strangers at any party, \emph{Amer.\ Math.\ Monthly} 66 (1959), 778--783.
\bibitem{Goo85} A.~W.~Goodman, Triangles in a complete chromatic graph with three colors, \emph{Discrete Math.} 57 (1985), 225--235.
\bibitem{GC24} P.~J.~Goulart and Y.~Chen, Clarabel: an interior-point solver for conic programs with quadratic objectives, preprint, arXiv:2405.12762, 2024.
\bibitem{GG55} R.~E.~Greenwood and A.~M.~Gleason, Combinatorial relations and chromatic graphs, \emph{Canad.\ J.\ Math.} 7 (1955), 1--7.
\bibitem{HW20} X.~He and Y.~Wigderson, Multicolor Ramsey numbers via pseudorandom graphs, \emph{Electron.\ J.\ Combin.} 27(1) (2020), \#P1.32.
\bibitem{HH18} Q.~Huangfu and J.~A.~J.~Hall, Parallelizing the dual revised simplex method, \emph{Math.\ Program.\ Comput.} 10 (2018), 119--142.
\bibitem{IMM18} A.~Ignatiev, A.~Morgado and J.~Marques-Silva, PySAT: a Python toolkit for prototyping with SAT oracles, in: \emph{Theory and Applications of Satisfiability Testing -- SAT 2018}, Lecture Notes in Comput.\ Sci.\ 10929, Springer, 2018, 428--437.
\bibitem{Kuz16} E.~Kuznetsov, Computational lower limits on small Ramsey numbers, arXiv:1505.07186v5, 2016.
\bibitem{Les} A.~Lesser, Theoretical and computational aspects of Ramsey theory, Examensarbeten i Matematik 3, Stockholms Universitet, 2001.
\bibitem{LP21} B.~Lidick\'y and F.~Pfender, Semidefinite programming and Ramsey numbers, \emph{SIAM J.\ Discrete Math.} 35 (2021), 2328--2344; arXiv:1704.03592.
\bibitem{Rad26} S.~P.~Radziszowski, Small Ramsey numbers, \emph{Electron.\ J.\ Combin.}, Dynamic Survey DS1.18 (revision 18, 24 April 2026), doi:10.37236/21.
\bibitem{Raz07} A.~A.~Razborov, Flag algebras, \emph{J.\ Symbolic Logic} 72 (2007), 1239--1282.
\end{thebibliography}

\end{document}
